% Options for packages loaded elsewhere
\PassOptionsToPackage{unicode}{hyperref}
\PassOptionsToPackage{hyphens}{url}
\PassOptionsToPackage{dvipsnames,svgnames,x11names}{xcolor}
\documentclass[
  10pt,
]{article}
\usepackage{xcolor}
\usepackage[margin=2.5cm]{geometry}
\usepackage{amsmath,amssymb}
\setcounter{secnumdepth}{-\maxdimen} % remove section numbering
\usepackage{iftex}
\ifPDFTeX
  \usepackage[T1]{fontenc}
  \usepackage[utf8]{inputenc}
  \usepackage{textcomp} % provide euro and other symbols
\else % if luatex or xetex
  \usepackage{unicode-math} % this also loads fontspec
  \defaultfontfeatures{Scale=MatchLowercase}
  \defaultfontfeatures[\rmfamily]{Ligatures=TeX,Scale=1}
\fi
\usepackage{lmodern}
\ifPDFTeX\else
  % xetex/luatex font selection
\fi
% Use upquote if available, for straight quotes in verbatim environments
\IfFileExists{upquote.sty}{\usepackage{upquote}}{}
\IfFileExists{microtype.sty}{% use microtype if available
  \usepackage[]{microtype}
  \UseMicrotypeSet[protrusion]{basicmath} % disable protrusion for tt fonts
}{}
\makeatletter
\@ifundefined{KOMAClassName}{% if non-KOMA class
  \IfFileExists{parskip.sty}{%
    \usepackage{parskip}
  }{% else
    \setlength{\parindent}{0pt}
    \setlength{\parskip}{6pt plus 2pt minus 1pt}}
}{% if KOMA class
  \KOMAoptions{parskip=half}}
\makeatother
\usepackage{longtable,booktabs,array}
\usepackage{caption}
\captionsetup[table]{skip=6pt}
\newcounter{none} % for unnumbered tables
\usepackage{calc} % for calculating minipage widths
% Correct order of tables after \paragraph or \subparagraph
\usepackage{etoolbox}
\makeatletter
\patchcmd\longtable{\par}{\if@noskipsec\mbox{}\fi\par}{}{}
\makeatother
\usepackage{graphicx}
\makeatletter
\newsavebox\pandoc@box
\newcommand*\pandocbounded[1]{% scales image to fit in text height/width
  \sbox\pandoc@box{#1}%
  \Gscale@div\@tempa{\textheight}{\dimexpr\ht\pandoc@box+\dp\pandoc@box\relax}%
  \Gscale@div\@tempb{\linewidth}{\wd\pandoc@box}%
  \ifdim\@tempb\p@<\@tempa\p@\let\@tempa\@tempb\fi% select the smaller of both
  \ifdim\@tempa\p@<\p@\scalebox{\@tempa}{\usebox\pandoc@box}%
  \else\usebox{\pandoc@box}%
  \fi%
}
% Set default figure placement to htbp
\def\fps@figure{htbp}
\makeatother
\setlength{\emergencystretch}{3em} % prevent overfull lines
\providecommand{\tightlist}{%
  \setlength{\itemsep}{0pt}\setlength{\parskip}{0pt}}
\usepackage{float}
\floatplacement{figure}{H}
\usepackage{caption}
\captionsetup{font=small}
\renewcommand{\arraystretch}{1.15}
\usepackage{bookmark}
\IfFileExists{xurl.sty}{\usepackage{xurl}}{} % add URL line breaks if available
\urlstyle{same}
% fallback for those not using the hyperref driver hyperxmp:
\makeatletter
\@ifundefined{xmpquote}{\newcommand{\xmpquote}[1]{#1}}{}
\makeatother
\hypersetup{
  pdftitle={Losing Leverage: A Game-Theoretic Simulation of Power After Full Automation},
  pdfauthor={Andrei Taranu (Eight Rice) (contact@autonet.computer, ORCID 0009-0008-2902-0845)},
  colorlinks=true,
  linkcolor={blue},
  filecolor={Maroon},
  citecolor={Blue},
  urlcolor={blue},
  pdfcreator={LaTeX via pandoc}}

\title{Losing Leverage: A Game-Theoretic Simulation of Power After Full Automation}
\author{Andrei Taranu (Eight Rice) (contact@autonet.computer, ORCID 0009-0008-2902-0845)}
\date{October 2026}

\begin{document}
\maketitle

\emph{Code, data and results: \url{https://github.com/autonet-code/leverage-sims}. All numbers in
this paper are reproducible from the code with the seeds given in Appendix D.}

\subsection{Abstract}\label{abstract}

Throughout history, rulers have needed large numbers of people: to work, to pay taxes and to enforce
orders. That need is what made them bargain with, feed and fear their populations. If AI and robots
can do nearly all productive and military work, the need may disappear. This paper asks what
protects people once it does.

We built a game-theoretic Monte Carlo simulation of six world blocs from 2026 to 2075. It tracks:

\begin{itemize}
\tightlist
\item
  AI capability and robot build-out;
\item
  the loss of the public's economic and military leverage;
\item
  the erosion of democracy;
\item
  purges inside ruling groups;
\item
  the choices those groups make once their populations are no longer needed.
\end{itemize}

It draws on historical base rates and includes four restraints on rulers found by a search of the
historical record with rules fixed in advance. We cross-checked it against a pre-specified
elicitation from a judgment model.

Like war games and the scenario side of climate research, the model cannot predict when events will
happen. It can compare futures and show which assumptions drive them. We report it that way.

\textbf{Main findings:}

\begin{enumerate}
\def\labelenumi{\arabic{enumi}.}
\tightlist
\item
  \textbf{Losing leverage removes the public's protection, and nothing structural replaces it.} By
  2075 the US or China public is disempowered in 94\% of runs, counting plain autocracy, and
  stripped of rights or worse in 73\%.
\item
  \textbf{The outcome then rests on a few people.} Once the public has no leverage, what decides
  whether rulers escalate to mass killing is their own restraint, incentives and internal politics,
  which no data measure.
\item
  \textbf{The result does not hinge on how closure is reached.} Building a dedicated automated loop
  instead of converting the existing economy brought closure about 2.5 years earlier but left the
  risk of mass killing almost unchanged. Faster AI progress does raise the risk.
\item
  \textbf{One intervention changes the incentives.} An economic network outside state and corporate
  control that pays its output to households makes it almost free for rulers to keep people alive.
  If it carries 50\% of US-bloc economic activity by the early 2030s, the risk of near-total
  depopulation falls by more than a quarter, and by about half for the period to 2050. Arriving in
  2040 loses about 40\% of that effect.
\end{enumerate}

\textbf{Where the assumptions lead.} The results hold under seventeen stated assumptions: eleven
that define the scenario and six that ground ruling-group behavior in evidence. Under them:

\begin{itemize}
\tightlist
\item
  people stop being needed in the leading bloc by a median year of 2034;
\item
  the probability of a deliberate loss of 10\% or more of the US or China bloc's population by 2075
  is about 0.44;
\item
  the probability that such a loss reaches 99.9\% or more (roughly 8 billion people to 8 million) is
  about 0.38 for the US or China bloc and 0.09 for the world.
\end{itemize}

The near-total level rests on one assumption: that rulers who have started killing fear the
survivors and finish the job. Without it, no run passes 99.9\%. These probabilities are far above
published expert and superforecaster estimates. The model's rate of lethal policy within 15 years of
closure is also 1.7 to 2.4 times the judgment model's. They show where the assumptions lead, not a
forecast.

\begin{center}\rule{0.5\linewidth}{0.5pt}\end{center}

\subsection{Key findings}\label{key-findings}

\emph{All findings below are model results under the assumptions in Section 4. The comparisons (what
changes the outcome, and by how much) are more robust than the absolute probabilities and dates.}

\begin{enumerate}
\def\labelenumi{\arabic{enumi}.}
\tightlist
\item
  \textbf{People lose their leverage in the early-to-mid 2030s.}

  \begin{itemize}
  \tightlist
  \item
    We call the point where a bloc's core supply chain needs less than 20\% of its 2026 human labor
    \textbf{closure}. The first bloc reaches it in a median year of 2034.0 (10th to 90th percentile
    2032.0 to 2039.25, among runs where closure occurs).
  \item
    Democracy erodes as leverage goes. In the US bloc, democracy breaks down by 2040 in 63\% of
    runs. Across all blocs, only 26\% of first post-closure decisions are taken while the deciding
    bloc is still a democracy.
  \end{itemize}
\item
  \textbf{No ruling group chooses depopulation first. It arises through escalation.}

  \begin{itemize}
  \tightlist
  \item
    First choices after closure are to keep the status quo, to serve the public, or to
    \textbf{warehouse} people (minimal provision and no role in the economy).
  \item
    Under unrest and fear of the survivors, 36\% of warehousing regimes later escalate to killing.
  \end{itemize}
\item
  \textbf{Near-total depopulation has a probability of more than one in three.} The probability of a
  deliberate loss of 99.9\% or more is:

  \begin{itemize}
  \tightlist
  \item
    0.378 for the US or China bloc by 2075;
  \item
    0.223 by 2050;
  \item
    0.607 for at least one of six blocs.
  \end{itemize}

  Without the threat-elimination assumption (A11), no run reaches 99.9\%. A deliberate loss of 10\%
  or more stays at a little more than two in five (0.44).
\item
  \textbf{Loss of standing is nearly universal.} By 2075, the US or China public is stripped of
  rights or worse in 73\% of runs, and disempowered, counting plain autocracy, in 94\%. Prolonged
  disempowerment is a catastrophe in its own right, not only a step toward mass killing.
\item
  \textbf{The outcome hinges on a few people.}

  \begin{itemize}
  \tightlist
  \item
    By 2075, the leader holds 95\% or more of coercive control in 58\% of US-or-China runs.
  \item
    Across plausible parameter values, the middle 95\% of estimates runs from 0.02 to 0.95. Given
    assumptions A4 to A6, which remove the public's leverage, what decides a run is the moral
    restraint, incentives and purge behavior of the ruling group.
  \end{itemize}
\item
  \textbf{A parallel economy based on human labor alone does not protect people.}

  \begin{itemize}
  \tightlist
  \item
    People trading among themselves outside the automated economy end as a tolerated community in
    about 0.2\% of runs.
  \item
    After closure, they have nothing the automated side needs.
  \end{itemize}
\item
  \textbf{An income network independent of the state lowers the risk substantially.}

  \begin{itemize}
  \tightlist
  \item
    It does so by paying households directly. It needs to be large by the early 2030s.
  \item
    20\% of economic activity buys about 70\% of the effect of 50\%, and even 5\% helps.
  \item
    It does not prevent crackdowns or change who controls force, and a probability of about 0.28
    remains.
  \item
    We do not assess whether such adoption is achievable. It is far above historical diffusion
    rates.
  \end{itemize}
\end{enumerate}

\begin{center}\rule{0.5\linewidth}{0.5pt}\end{center}

\subsection{Disclosures}\label{disclosures}

\textbf{Competing interests.} The author develops Autonet, a decentralized AI project of the general
kind tested in Section 6.2. The network in Section 6.2 is modeled as a generic class of system
defined only by its effects, not as Autonet or any other specific product. The author has no
affiliation with TypeSafe, the maker of the Jev judgment model.

\textbf{Use of AI tools.} AI tools assisted with coding, and source checking; the author is
responsible for all content.

\begin{center}\rule{0.5\linewidth}{0.5pt}\end{center}

\subsection{1 Introduction}\label{introduction}

\subsubsection{1.1 The scenario}\label{the-scenario}

The hypothesis we test is a chain of events, each step ordinary on its own:

\begin{enumerate}
\def\labelenumi{\arabic{enumi}.}
\tightlist
\item
  \textbf{Firms automate because their competitors do.} Profits concentrate in a few AI providers.
\item
  \textbf{Broad economic participation collapses.} As people lose income, consumer businesses lose
  customers. Business-to-business firms follow a little later.
\item
  \textbf{Displaced people protest, then revolt.} They cannot threaten a state whose force is
  automated and whose tax base is the AI firms.
\item
  \textbf{Rulers ask why they should keep paying for people they do not need.} People become mouths
  to feed, and loud ones. Depopulation can be justified internally as an ecological reset or as
  kindness.
\item
  \textbf{Partial measures leave survivors who remember.} Survivors are witnesses and potential
  avengers, inside a bloc and across borders. This logic pushes toward near-total reduction, with a
  small captive remnant kept without rights.
\end{enumerate}

We ask three questions. How likely is this chain under explicit assumptions? When would it happen?
Which links drive it?

\subsubsection{1.2 Why this is worth modeling}\label{why-this-is-worth-modeling}

Most quantitative work on catastrophic AI risk studies misaligned AI: systems that pursue goals
their builders did not intend (Ord 2020; Carlsmith 2022). Economic work studies job loss and growth.
The mechanism here is different. The AI works as intended, and humans direct it. The danger comes
from the people at the top no longer needing everyone else.

This mechanism has been described qualitatively:

\begin{itemize}
\tightlist
\item
  \emph{The Intelligence Curse} (Drago and Laine 2025);
\item
  \emph{Gradual Disempowerment} (Kulveit et al.~2025);
\item
  \emph{AI-Enabled Coups} (Davidson, Finnveden and Hadshar 2025).
\end{itemize}

We found no published quantitative model of it. Empirical work so far measures milder effects. For
example, AI surveillance in China helps the state suppress protest (Beraja et al.~2023). This paper
is a first attempt at a full quantitative model.

\subsubsection{1.3 How to read the numbers}\label{how-to-read-the-numbers}

Every probability in this paper is conditional on the assumptions in Section 4. Each is stated
plainly, so a reader can reject one and see roughly what follows. Section 5.8 and Appendix E.5
remove threat elimination (A11), each of the six behavioral assumptions (B1 to B6), the environment
trends (A7), the means of mass killing (A8), collusion (A9) and several mechanisms. They also revert
A1, A2, A3 and A6 to earlier, milder forms. A4 and A5 are not tested by removal, for these reasons:

\begin{itemize}
\tightlist
\item
  \textbf{A4.} Once force is automated, refusal by the people who carry out orders no longer
  applies, and nearly every path to the outcome requires closure (Section 5.3). Keeping humans in
  the loop of automated force is a policy lever (Section 6.1).
\item
  \textbf{A5.} Keeping a safety net is also a policy choice, and it is not current reality.
\end{itemize}

Section 8.5 (item 3) reports the related model checks.

\begin{center}\rule{0.5\linewidth}{0.5pt}\end{center}

\subsection{2 The core argument: why people have power, and how automation removes
it}\label{the-core-argument-why-people-have-power-and-how-automation-removes-it}

Rulers have always depended on large numbers of people. They needed farmers to grow food, workers to
make goods, taxpayers to fund the state, and soldiers and police to enforce orders. That dependence
gave ordinary people leverage in three forms.

\begin{enumerate}
\def\labelenumi{\arabic{enumi}.}
\tightlist
\item
  \textbf{Labor leverage.} People can stop working. Strikes, slowdowns and emigration impose costs
  that rulers must weigh.
\item
  \textbf{Revolt leverage.} People can fight. A regime that faces a credible uprising must bargain
  or repress, and repression needs enforcers.
\item
  \textbf{Enforcer leverage.} The enforcers are people too. Soldiers can refuse to fire, and police
  can defect. Most dictators who lose power by unconstitutional means are removed by insiders: 205
  of 316 between 1946 and 2008 (Svolik 2012). Campaigns facing a violent crackdown were more than
  four times as likely to succeed when security forces defected (Stephan and Chenoweth 2008).
\end{enumerate}

Democracy, welfare and the rule of law grew where these forms of leverage were strong. They were not
gifts but bargains that rulers accepted because the alternative cost more. This is the core claim of
several literatures:

\begin{itemize}
\tightlist
\item
  selectorate theory (Bueno de Mesquita et al.~2003);
\item
  economic theories of democratization (Acemoglu and Robinson 2006);
\item
  work on how taxation built modern states (Tilly 1990).
\end{itemize}

Under assumptions A1 to A4, automation removes all three forms of leverage:

\begin{itemize}
\tightlist
\item
  When AI and robots run the core supply chain (mining, energy, chips, manufacturing, logistics,
  construction, maintenance), labor leverage goes.
\item
  When security is automated, revolt leverage goes. A crowd cannot overpower a fleet of drones, and
  drones do not need to be persuaded.
\item
  When no human operator stands behind each weapon, enforcer leverage goes. The only people left who
  could refuse are the handful who authorize operations at the top.
\end{itemize}

\textbf{The key structural point.} None of this requires rulers to be unusually cruel. It requires
only that the costs and benefits they face change.

\begin{itemize}
\tightlist
\item
  \textbf{Today}, keeping a population fed and housed is cheaper than the alternative, and the
  population is needed anyway.
\item
  \textbf{After closure}, keeping people costs resources, land and the risk of unrest, and returns
  nothing the rulers need.
\end{itemize}

What happens then depends on the restraint of the people in charge. The historical record of people
who have held unchecked power gives little reason for comfort (Section 4.2).

\begin{center}\rule{0.5\linewidth}{0.5pt}\end{center}

\subsection{3 Method}\label{method}

\subsubsection{3.1 Three lines of evidence}\label{three-lines-of-evidence}

\begin{enumerate}
\def\labelenumi{\arabic{enumi}.}
\item
  \textbf{Historical and empirical base rates.} We used the best available data for each link:

  {\def\LTcaptype{none} % do not increment counter
  \begin{longtable}[]{@{}
    >{\raggedright\arraybackslash}p{(\linewidth - 2\tabcolsep) * \real{0.5000}}
    >{\raggedright\arraybackslash}p{(\linewidth - 2\tabcolsep) * \real{0.5000}}@{}}
  \toprule\noalign{}
  \begin{minipage}[b]{\linewidth}\raggedright
  Domain
  \end{minipage} & \begin{minipage}[b]{\linewidth}\raggedright
  Source
  \end{minipage} \\
  \midrule\noalign{}
  \endhead
  \bottomrule\noalign{}
  \endlastfoot
  AI capability trends & METR task-horizon measurements (Kwa et al.~2025) \\
  Robot production & IFR (2025) \\
  Compute shares & Epoch AI (2025) \\
  Military spending & SIPRI (2025) \\
  Democratic erosion & V-Dem (Nord et al.~2026) \\
  Obedience and refusal & Milgram (1974) and replications (Haslam, Loughnan and Perry 2014; Burger
  2009); Browning (1992) \\
  Elite purges & Sudduth (2017), Goldring \\
  Mass killing & Harff (2003), Valentino, Huth and Balch-Lindsay (2004), Straus (2006) \\
  \end{longtable}
  }

  Current policy trajectories, as of September 2026, set the starting conditions.
\item
  \textbf{A game-theoretic Monte Carlo simulation.} Six blocs evolve year by year. In each bloc, a
  ruling group of individually simulated members (not modeled on real people) weighs:

  \begin{itemize}
  \tightlist
  \item
    the value of freed land and resources;
  \item
    their own moral cost;
  \item
    the threat of revolt;
  \item
    outside pressure that someone can actually enforce;
  \item
    the power of rival blocs.
  \end{itemize}

  Grievance, insurgency, democratic erosion, purges, war and cross-border campaigns arise inside the
  model rather than being imposed.
\item
  \textbf{A judgment model as a cross-check.} Jev (TypeSafe) is a language model designed to give
  calibrated probability judgments. We fixed the design before sending any query and committed it to
  the repository: 648 scenario cells, each asked in 3 paraphrases, plus a 648-cell control arm in
  which humans still do 80\% of the work, for 3,888 queries in total.

  We did not calibrate the simulation's outputs to Jev. Jev does inform two inputs:

  \begin{itemize}
  \tightlist
  \item
    the share of elites who refuse to order killing, mixed with historical priors;
  \item
    the effect of provision on revolt, as a 0.6 / 0.4 mixture of historical evidence and Jev.
  \end{itemize}

  Everywhere else, Jev serves only as a check, mainly on which factors raise or lower risk.
\end{enumerate}

\subsubsection{3.2 The simulation in brief}\label{the-simulation-in-brief}

\textbf{Blocs} (populations from UN DESA 2024).

{\def\LTcaptype{none} % do not increment counter
\begin{longtable}[]{@{}lll@{}}
\toprule\noalign{}
Bloc & Members & Population 2026 \\
\midrule\noalign{}
\endhead
\bottomrule\noalign{}
\endlastfoot
US bloc & US, Japan, South Korea, Taiwan & 0.53 bn \\
China bloc & China & 1.42 bn \\
Europe plus & EU, UK, Canada, Australia, New Zealand, Switzerland, Norway & 0.62 bn \\
Russia and MENA & Russia, Middle East, North Africa & 0.99 bn \\
South Asia & & 1.97 bn \\
Global South & Remaining countries & 2.67 bn \\
\end{longtable}
}

Each bloc has these parts:

\begin{enumerate}
\def\labelenumi{\arabic{enumi}.}
\tightlist
\item
  \textbf{An economy of nine segments}: mining, energy, chips, manufacturing, logistics,
  construction, maintenance, security and administration. Each segment's automated share grows with
  AI capability and robot supply. The limits on that growth are:

  \begin{itemize}
  \tightlist
  \item
    robots building robots;
  \item
    energy;
  \item
    chip fabrication;
  \item
    mining;
  \item
    the bloc's share of world input supply.
  \end{itemize}
\item
  \textbf{A human-dependence index.} This is the share of the core supply chain's 2026 human labor
  that is still needed, counting labor embedded in imports. Closure is the point where it falls
  below 20\%. Section 4.4 describes how a dedicated automated loop enters this index.
\item
  \textbf{A public.} Displacement and cuts to provision raise grievance. Grievance raises hostility,
  and hostility raises the yearly chance of insurgency (we call such yearly rates \textbf{hazards}).
  Insurgency drives repression.
\item
  \textbf{A democracy score} on the V-Dem liberal democracy scale. It falls as the public's labor
  and revolt leverage fall. We count a score below 0.30 as a breakdown, roughly Hungary in 2024 or
  Turkey in 2015.
\item
  \textbf{A ruling coalition.} In autocracies it has up to 21 members; in democracies, a wider elite
  of 51 decides through veto players. Each member has:

  \begin{itemize}
  \tightlist
  \item
    a moral cost of ordering killing;
  \item
    a degree of concern for the public;
  \item
    a discount rate;
  \item
    a share of control over coercive force.
  \end{itemize}

  Members can purge each other, and purges can trigger counter-coups. A leader who holds half or
  more of coercive control rules alone. The inner circle can still block a decision, but its chance
  of success falls as automation grows.
\item
  \textbf{Military power.} This depends on military spending, economic output and the automated
  share of security.
\end{enumerate}

\textbf{Choices.} After closure, each ruling coalition repeatedly chooses among six options, from
mildest to harshest:

\begin{enumerate}
\def\labelenumi{\arabic{enumi}.}
\tightlist
\item
  serve the public;
\item
  live off rents;
\item
  keep the status quo;
\item
  warehouse people;
\item
  neglect them;
\item
  depopulate.
\end{enumerate}

How a regime decides depends on its type:

\begin{itemize}
\tightlist
\item
  \textbf{Democracies} need near-unanimity among 3 to 5 veto players (Tsebelis 2002). Their number
  shrinks as democracy erodes.
\item
  \textbf{Oligarchies} need a 50 to 67\% supermajority for a harsher move, with each member's vote
  weighted by their control of force.
\item
  \textbf{A personalist leader} decides alone, subject to the inner-circle veto.
\end{itemize}

\textbf{Endgame.} The simulation tracks deaths from five channels:

\begin{enumerate}
\def\labelenumi{\arabic{enumi}.}
\tightlist
\item
  executed depopulation orders;
\item
  economic strangulation and neglect;
\item
  an abstract means of rapid mass killing, whose availability rises with AI capability (no
  operational detail is modeled);
\item
  war, including war between coalitions that have both chosen depopulation, up to nuclear exchange;
\item
  cross-border campaigns against blocs that cannot resist.
\end{enumerate}

Only three things can stop a campaign: a successful revolt, a repelled attack, or nuclear
retaliation.

\textbf{Scenario tree.} The bloc simulation sits inside a broader scenario tree, which adds three
things:

\begin{itemize}
\tightlist
\item
  \textbf{Earlier stages.} These are the automation race, the collapse of broad economic
  participation, and broad repression by regimes still staffed by humans.
\item
  \textbf{Background catastrophes.} AI takeover not directed by rulers and non-AI global catastrophe
  are competing end points. A run that hits one first is not counted toward depopulation.
\item
  \textbf{A counterfactual US history.} If the US first becomes a broadly repressive regime, the
  model continues from a world in which the US is autocratic at closure, for all six blocs.
\end{itemize}

\subsubsection{3.3 Sample size and what the intervals
mean}\label{sample-size-and-what-the-intervals-mean}

The main run draws 3,000 sets of uncertain parameter values (we call each set a \textbf{draw}). Each
draw is run for 200 paths spread over 8 random seeds, giving 3,000 x 200 = 600,000 simulated futures
from October 2026 to the end of 2075. We report three kinds of interval:

\begin{enumerate}
\def\labelenumi{\arabic{enumi}.}
\tightlist
\item
  \textbf{Simulation error} (bootstrap 95\% interval). This shows how precisely the sample pins down
  the average. For the headline it is 0.366 to 0.391.
\item
  \textbf{Parameter spread} (95\% range of per-draw probabilities). This shows how much the answer
  depends on which plausible parameter values are true. For the headline it is 0.021 to 0.951.
\item
  \textbf{Structural range.} This is the spread across 46 alternative model versions. It runs from
  0.000 (threat elimination switched off) to 0.875 (every pessimistic option switched on).
\end{enumerate}

\textbf{What the precision means.} The small simulation error says nothing about accuracy. The
headline is an average over judgment priors, and shifting those priors shifts it. Per-draw values
are close to all-or-nothing: a draw's parameters either permit an executed order or they do not. We
therefore quote ``about 0.38'' in the text and keep three decimals for tables.

\subsubsection{3.4 Evidence labels}\label{evidence-labels}

Every parameter is labeled by the kind of evidence behind it:

\begin{itemize}
\tightlist
\item
  measured;
\item
  strong analog;
\item
  weak analog;
\item
  judgment.
\end{itemize}

Where a situation has no historical precedent, the parameter gets a wide range. It is never pushed
toward zero because the situation is new. The whole scenario is unprecedented: populations have
never lost both their economic and their military relevance.

Gaps are filled with documented human tendencies:

\begin{itemize}
\tightlist
\item
  in-group and out-group behavior;
\item
  dehumanization and moral disengagement (Bandura 1999);
\item
  elite self-interest;
\item
  the over-representation of people with weak moral restraint among those who hold unchecked power.
\end{itemize}

Appendix A lists the key parameters with their labels. Supplement S1 lists all of them.

\subsubsection{3.5 Verification}\label{verification}

\begin{itemize}
\tightlist
\item
  \textbf{Independent audit.} An audit reran the full pipeline of the version without brakes and
  reproduced its headline numbers exactly.
\item
  \textbf{Backward compatibility.} With the brakes switched off, the code reproduces the version
  without brakes exactly. Switching off B1 to B6 as well reproduces the version before them exactly.
  With the three structural changes of Section 4.4 switched off, the code reproduces the brakes-only
  version exactly (0.36744 and 0.08784).
\item
  \textbf{CPU and GPU.} The GPU and CPU implementations agree within simulation error.
\item
  \textbf{Bookkeeping.} Outcome categories sum to one, loss thresholds are nested correctly, and
  deaths are not double counted across blocs.
\end{itemize}

The bookkeeping bugs found by the final two audits are listed in Appendix C. None changed the
headline by more than 0.002. Earlier design errors, which did change results, are listed there too.

\begin{center}\rule{0.5\linewidth}{0.5pt}\end{center}

\subsection{4 Assumptions}\label{assumptions}

\subsubsection{4.1 Scenario assumptions}\label{scenario-assumptions}

These eleven assumptions define the scenario. They were fixed before the final runs.

\textbf{A1. AI capability keeps growing exponentially, and the growth can speed up.}

\begin{itemize}
\tightlist
\item
  The length of task AI can complete on its own has doubled about every four months since 2023.
  Averaged over 2019 to 2025, it doubled about every seven months (METR; Kwa et al.~2025).
\item
  The doubling time can shrink further as AI speeds up AI research.
\item
  In the model, faster AI progress can only bring closure earlier, never later.
\end{itemize}

\textbf{A2. Physical automation is a top priority in the US and China from 2026.} Robot output grows
exponentially. It is limited only by how fast robots can build robots and by documented physical
limits.

\textbf{A3. Security automation needs no separate investment and gets priority.} Military and police
automation has no separate, slower curve. It follows from general capability and robot supply, and
security gets first claim on that supply.

\textbf{A4. Refusal by soldiers and police stops mattering once autonomous systems are fielded at
scale.}

\begin{itemize}
\tightlist
\item
  No human operator stands behind each drone. Autonomous targeting is already used in Ukraine.
\item
  The only refusal that still counts is among the few people who authorize operations at the top.
\item
  Repression aimed at reducing numbers needs no checkpoints or detention, and therefore no human
  staff.
\item
  In the model, switching off defection by security forces leaves the headline almost unchanged
  (0.378 against 0.384 at the same sample size), because security is largely automated before any
  decision is taken.
\end{itemize}

\textbf{A5. There is no safety net by default.} Provision starts on the 2026 trajectory of cuts. In
the US, the 2025 reconciliation act cuts federal Medicaid spending by about \$911 billion and food
assistance by about \$186 billion over ten years. The Congressional Budget Office expects 10 million
more people to be uninsured by 2034 (CBO 2025). Keeping people fed is a choice rulers must actively
make, and cheaper production does not imply willingness to provide.

\textbf{A6. Power inside a ruling group equals control of automated force.} Money and resources
based on human labor count for nothing once labor is not needed. Democracy survives only while the
public holds real leverage.

\textbf{A7. The world environment is deteriorating.} Current conditions set the model's trends for
extremism, war, emergency powers and moral restraint:

\begin{itemize}
\tightlist
\item
  State-based armed conflicts reached 61 in 2024, the most since 1946 (Rustad 2025).
\item
  Global peacefulness is at its lowest since the Global Peace Index began (IEP 2026).
\item
  The US withdrew from 66 international organizations in January 2026 (White House 2026).
\item
  In 2023, 23\% of Americans agreed that ``true American patriots may have to resort to violence to
  save our country'', up from 15\% in 2021 (PRRI 2023).
\item
  A record-strength El Niño is forecast for late 2026 (NOAA 2026).
\item
  Far-right parties keep gaining. For example, the AfD won 43.8\% in the September 2026
  Saxony-Anhalt state election.
\end{itemize}

\textbf{A8. The headline is near-total depopulation: a loss of 99.9\% or more, roughly 8 billion to
8 million.} We also report losses of 10\%, 50\% and 90\%. The means of rapid mass killing is a
single abstract availability parameter.

\textbf{A9. Mechanisms stack.}

\begin{itemize}
\tightlist
\item
  Once a population is targeted, economic strangulation, neglect, the abstract means and war can all
  act at the same time.
\item
  Rulers of two blocs that have both chosen depopulation may tacitly tolerate attacks on each
  other's populations, up to nuclear exchange.
\item
  Rulers weigh the risk to their own lives as a large cost rather than an absolute limit, because
  they can shelter in protected bunkers.
\end{itemize}

\textbf{A10. ``Unprecedented'' means uncertain, not unlikely.} See Section 3.4.

\textbf{A11. Threat elimination.}

\begin{itemize}
\tightlist
\item
  Survivors are witnesses and potential avengers, so partial depopulation tends to escalate toward
  completion. Perpetrators of mass killing often target the people they see as future threats
  (Valentino, Huth and Balch-Lindsay 2004; Straus 2006). How strongly this scales to whole
  populations is judgment.
\item
  The fear of retribution grows with the scale of what has been done.
\item
  A small remnant (0.03 to 0.1\% of the population) may be kept, without rights. We count the
  remnant as disempowered, not as survivors.
\end{itemize}

\subsubsection{4.2 Behavioral assumptions}\label{behavioral-assumptions}

These six assumptions ground the behavior of ruling groups in evidence. They were adopted before the
brake search in Section 4.3.

\textbf{B1. The moral restraint of people at the top is anchored in how autocrats actually behave,
not in how average people would.}

\begin{itemize}
\tightlist
\item
  In autocracies, the leader's moral cost of ordering killing is scaled down by a factor drawn from
  0.1 to 0.6. Other coalition members are scaled down less. Democratic leaders are unchanged.
\item
  The evidence supports the direction, but it is thin:

  \begin{itemize}
  \tightlist
  \item
    Psychopathy affects about 1.2\% of the general population (Sanz-García et al.~2021).
    Psychopathic traits appeared more often in a non-random sample of 203 corporate managers
    (Babiak, Neumann and Hare 2010). The link to emerging as a leader is weak overall (Landay, Harms
    and Credé 2019).
  \item
    A synthesis of tyrants' psychology finds paranoia that worsens once power is secured (Glad
    2002).
  \item
    Russia's leadership has kept its course through more than a million killed and wounded (UK
    Ministry of Defence 2025).
  \item
    North Korea runs political prison camps that a UN inquiry classed as crimes against humanity (UN
    Commission of Inquiry 2014).
  \item
    Stalin's purges executed 681,692 people in 1937 and 1938 by official count (Getty, Rittersporn
    and Zemskov 1993).
  \end{itemize}
\item
  No evidence measures any leader's moral cost of killing an entire population. The size of the
  effect is judgment.
\end{itemize}

\textbf{B2. Concentration of power emerges from incentives.}

\begin{itemize}
\tightlist
\item
  Each coalition member's weight is their share of control over coercive force.
\item
  Purges happen at a rate that rises as coercion stops depending on human networks. Dictators purge
  when the elite's capacity to oust them is weak (Sudduth 2017), and they trade competence for
  loyalty (Egorov and Sonin 2011).
\item
  Purges can trigger counter-coups.
\item
  Nothing forces a single ruler. Whether one emerges is an outcome of the model.
\end{itemize}

\textbf{B3. Outside pressure counts only if someone can enforce it.} International condemnation
matters in proportion to the condemning party's military power relative to the target. It fades to
nothing when one bloc dominates.

\textbf{B4. Fear of retribution grows with the scale of what has been done.} The more people a
regime has killed, the more it fears the survivors, and the more it purges its own ranks.

\textbf{B5. The fear of survivors extends abroad.} A depopulating coalition may attack another bloc
whose population it sees as a future threat, under two conditions:

\begin{itemize}
\tightlist
\item
  its share of the two blocs' combined military power exceeds a threshold drawn between 70\% and
  95\%;
\item
  the expected gain outweighs the risk of nuclear retaliation.
\end{itemize}

Each bloc's ability to resist depends on its own automated military power. Personalist and military
regimes start international conflicts more often than other regimes (Weeks 2012).

\textbf{B6. Physical bottlenecks are binding, and blocs differ.}

\begin{itemize}
\tightlist
\item
  Each bloc draws on a world pool of robot inputs according to its production and imports, and
  exporters restrict sales.
\item
  Industrial base sets how fast a bloc can bring production home.
\item
  Outside the leading blocs, AI capability lags the frontier by up to three years and robot
  capability by up to four.
\end{itemize}

\subsubsection{4.3 Brakes we searched for}\label{brakes-we-searched-for}

Every major change during model development removed a brake that contradicted conditions observed in
2026. To check the model for the opposite bias, we searched for restraining mechanisms it might be
missing. Before the search began, we fixed the rules for what would count
(\texttt{paper/brake\_search\_protocol.md}).

\textbf{Inclusion rules.} A candidate brake had to meet all three:

\begin{enumerate}
\def\labelenumi{\arabic{enumi}.}
\tightlist
\item
  \textbf{A documented mechanism.} A published causal account of how it restrains people who hold
  power.
\item
  \textbf{Two historical cases or more.} In each, the brake restrained rulers who had the capacity
  to carry out mass repression or killing. The case evidence has to credit the restraint to this
  mechanism.
\item
  \textbf{Model fit.} It can be expressed in the model, is not already carried by an existing
  parameter, and still operates once the public has lost its leverage (A3 to A6).
\end{enumerate}

\textbf{Process.}

\begin{itemize}
\tightlist
\item
  Each candidate was checked by two independent reviewers. One tested the mechanism and the cases.
  The other tested model fit. A candidate was included only if neither refuted it.
\item
  Strengths came from the cases. Where the cases could not bound a strength, it came from a Jev
  elicitation whose wording was fixed in advance, asked in three paraphrases.
\item
  No one chose brakes or strengths after the rules were fixed.
\end{itemize}

\textbf{Results.} The first round considered 15 candidates, and three passed. A review of the
strongest one found that it counted only one side of a coalition member's risk. A follow-up round
therefore searched for the opposite mechanism under the same rules, along with two related
candidates.

The four mechanisms included:

\begin{enumerate}
\def\labelenumi{\arabic{enumi}.}
\tightlist
\item
  \textbf{Self-risk (``I could be next'').}

  \begin{itemize}
  \tightlist
  \item
    \textbf{Mechanism.} Members of a ruling coalition expect that mass killing will make the leader
    more paranoid and the purges faster. They count the rise in their own risk as a cost when they
    vote. Members who expect to be exempt do not, and a shared depopulation ideology makes more
    members expect exemption.
  \item
    \textbf{Cases.} Thermidor in 1794 and the fall of Beria in 1953. Both involved elites already in
    danger themselves, so the brake's strength was set lower and wider than the first estimate.
  \end{itemize}
\item
  \textbf{Dissent risk.}

  \begin{itemize}
  \tightlist
  \item
    \textbf{Mechanism.} The mirror image of self-risk: members who object are themselves targeted,
    so others conform.
  \item
    \textbf{Cases.} Peng Dehuai's purge after the Lushan Conference in 1959, after which officials
    stayed silent through the worst famine years. Grigory Kaminsky and Osip Piatnitsky, arrested and
    shot after objecting to the 1937 purges. Hou Yuon, killed after opposing the Khmer Rouge
    evacuation of Phnom Penh.
  \item
    \textbf{Strength.} An objector's purge risk is set at about four times the normal rate, with a
    range of 1 to 30.
  \end{itemize}
\item
  \textbf{Leader mortality.} Leaders die or become incapacitated at 1 to 5\% a year at closure,
  falling as medicine improves. A successor may reverse course.
\item
  \textbf{Protective doctrine.} A ruling group may sincerely adopt a doctrine that makes killing far
  costlier, as Gorbachev's did in 1989. Jev puts the chance at about 6.5\% per nine years.
\end{enumerate}

\textbf{Rejected.} Fourteen candidates failed. Most failed because no historical case showed the
brake stopping rulers who could act, or because the model already carries it. They include:

\begin{itemize}
\tightlist
\item
  rulers' demand for an audience, servants or company (covered by the captive remnant);
\item
  keeping a population as a backup workforce;
\item
  church vetoes;
\item
  protecting one's own ethnic group;
\item
  fear of losing control of one's own AI;
\item
  AI systems refusing;
\item
  trade reputation;
\item
  fear of judgment after death, and the idea that a longer life lengthens the threat horizon.
\end{itemize}

Every candidate, with the reviewers' reasoning, is in \texttt{paper/brake\_search\_results.json} and
\texttt{paper/brake\_search\_followup.json}.

\textbf{Effect.} All four brakes are part of the main specification reported in this paper. Section
5.8 reports how much each one moves the results, alone and together.

\subsubsection{4.4 Structural changes after the brake
search}\label{structural-changes-after-the-brake-search}

A review of how closure was modelled, made before publication, led to three changes. All three are
part of the main specification. Table 4c (Section 5.8) shows the effect of each.

\textbf{1. A dedicated automated core loop.}

\begin{itemize}
\tightlist
\item
  \textbf{The earlier picture.} Earlier versions modelled closure as converting the existing
  economy. The automated share of each core segment grew from today's share, limited by how fast
  existing plants and equipment can be replaced or refitted.
\item
  \textbf{Why it was wrong.} A ruling group does not need to convert the whole economy to stop
  depending on human labor. It can build a new automated supply chain beside the existing one, sized
  only to keep itself supplied.
\item
  \textbf{How it is modelled.} The loop starts at zero: today's automated plants do not count. It
  takes a share of each bloc's new mining, energy and chip capacity, and it builds factories,
  logistics, construction, maintenance and administration in parallel. Investment, robot supply and
  what AI and robots can do at the time limit the build. Security automation stays shared with the
  rest of the economy (A3). A bloc's dependence on human labor is the lower of the two paths,
  conversion and loop.
\item
  \textbf{Loop size.} The loop is sized by the existing prior for the minimal self-sufficient core
  economy (q\_core): lognormal with a median of 0.15 of the full chain and a log-scale spread of 0.6
  (Appendix A.2). The change adds no new random draws.
\end{itemize}

\textbf{2. Own-industry compute feedback.}

\begin{itemize}
\tightlist
\item
  \textbf{The earlier picture.} All blocs followed the same capability trend, with fixed lags
  outside the leading blocs (B6). A bloc's own chip and energy build-out could not speed up its AI.
\item
  \textbf{How it is modelled.} When a bloc's own compute capacity (the lower of its chip and energy
  capacity) grows faster than the growth already built into the shared trend, the bloc gains extra
  capability doublings. It gains eta\_cf extra doublings per extra doubling of compute, with eta\_cf
  drawn from U(0.5, 1.5) (label C: the task horizon has doubled every 4 to 7 months while frontier
  training compute has grown about 4 to 5 times a year; Epoch AI 2024). The gain accrues only once
  the bloc's chip and energy sectors are mostly automated. Each year, other blocs close a share
  s\_sp of their gap to the leader, with s\_sp drawn from U(0.1, 0.5) (label J: theft, open weights
  and the movement of talent), and less when the leader is racing. The extra capability feeds the
  bloc's AI and robot capability, its core loop, its security automation and its military power. It
  has no ceiling of its own, so a bloc with enough compute can push past the hard limit on AI
  cognition sampled in some runs.
\item
  \textbf{Result: no pull-away between the US and China.} The extra capability is large: in the main
  run, each of the two blocs has a median of about 2.3 extra doublings in 2040. But both gain it at
  about the same time. The two close a median of 0.25 years apart, and in 2040 the leading bloc's
  extra capability exceeds the other's by a median of 0.06 doublings (by more than one doubling in
  0.6\% of runs).
\end{itemize}

\textbf{3. Wider physical floors under automation.}

\begin{itemize}
\tightlist
\item
  \textbf{The earlier picture.} Several parameters set the fastest pace at which automated industry
  can grow. Their lower tails were narrow, and the chip-fab doubling time had a hard floor at 0.5
  years.
\item
  \textbf{Why it was changed.} These are judgment parameters for a situation without precedent.
  Under A10, their ranges should be wide.
\item
  \textbf{The change} (label J, A10). The lower tails were widened. The 90th percentiles are
  unchanged.
\end{itemize}

{\def\LTcaptype{none} % do not increment counter
\begin{longtable}[]{@{}
  >{\raggedright\arraybackslash}p{(\linewidth - 6\tabcolsep) * \real{0.2500}}
  >{\raggedright\arraybackslash}p{(\linewidth - 6\tabcolsep) * \real{0.2500}}
  >{\raggedright\arraybackslash}p{(\linewidth - 6\tabcolsep) * \real{0.2500}}
  >{\raggedright\arraybackslash}p{(\linewidth - 6\tabcolsep) * \real{0.2500}}@{}}
\toprule\noalign{}
\begin{minipage}[b]{\linewidth}\raggedright
Parameter
\end{minipage} & \begin{minipage}[b]{\linewidth}\raggedright
Meaning
\end{minipage} & \begin{minipage}[b]{\linewidth}\raggedright
Old median (5th percentile)
\end{minipage} & \begin{minipage}[b]{\linewidth}\raggedright
New median (5th percentile)
\end{minipage} \\
\midrule\noalign{}
\endhead
\bottomrule\noalign{}
\endlastfoot
Td\_mine\_auto, Td\_en\_auto & Doubling time of automated mining and energy capacity, years & 1.5
(0.72) & 0.95 (0.25) \\
Td\_auto & Robot-stock doubling time once AI directs production, years & 1.0 (0.37) & 0.84 (0.25) \\
fact\_build\_a & Factory build time once the resource chain is automated, years & 0.50 (0.24) & 0.34
(0.10) \\
Td\_fab\_auto & Doubling time of automated chip-fab capacity, years & 1.04 (0.50, a hard floor) &
0.94 (0.25) \\
\end{longtable}
}

\textbf{Effect.} Together, the three changes bring closure about 2.5 years earlier. The headline
barely changes (Section 5.8, Table 4c).

\begin{center}\rule{0.5\linewidth}{0.5pt}\end{center}

\subsection{5 Results}\label{results}

\subsubsection{5.1 When people lose their leverage}\label{when-people-lose-their-leverage}

\textbf{Milestones} (in the bloc simulation run on its own):

{\def\LTcaptype{none} % do not increment counter
\begin{longtable}[]{@{}
  >{\raggedright\arraybackslash}p{(\linewidth - 4\tabcolsep) * \real{0.3333}}
  >{\raggedright\arraybackslash}p{(\linewidth - 4\tabcolsep) * \real{0.3333}}
  >{\raggedright\arraybackslash}p{(\linewidth - 4\tabcolsep) * \real{0.3333}}@{}}
\toprule\noalign{}
\begin{minipage}[b]{\linewidth}\raggedright
Milestone
\end{minipage} & \begin{minipage}[b]{\linewidth}\raggedright
Median year
\end{minipage} & \begin{minipage}[b]{\linewidth}\raggedright
10th to 90th percentile
\end{minipage} \\
\midrule\noalign{}
\endhead
\bottomrule\noalign{}
\endlastfoot
AI does half of software tasks end to end & 2028.0 & 2027.5 to 2029.0 \\
AI does half of all desk tasks & 2029.0 & 2028.0 to 2031.0 \\
AI reliably finishes a one-year project alone & 2029.5 & 2028.5 to 2031.5 \\
A general robot does half of physical tasks at human speed and cost & 2031.75 & 2030.0 to 2036.25 \\
Leading bloc's core chain needs less than half its 2026 labor & 2031.5 & 2030.25 to 2034.5 \\
\textbf{Closure: core chain needs less than 20\%} & \textbf{2034.0} & \textbf{2032.0 to 2039.25} \\
\end{longtable}
}

\begin{itemize}
\tightlist
\item
  \textbf{Closure in any bloc} (full scenario tree): 0.645 by 2035, 0.815 by 2040, and 0.878 by
  2075. Most runs without closure are runs in which the model samples a hard limit on robot
  dexterity or AI cognition.
\item
  \textbf{Which bloc closes first.} When closure happens, China closes first in 60\% of runs. This
  comes from the judgment prior that China's physical deployment lags the frontier by 0 years and
  the US by 0 to 1 year. That judgment is anchored on China's 54\% share of 2024 industrial robot
  installations (IFR 2025). The ordering is sensitive to it. The US and China close within a year of
  each other in 98\% of runs.
\item
  \textbf{Robot output.} Median production of general-purpose robots reaches about 0.6 million per
  year in 2030, 23 million in 2035 and 395 million in 2040.
\item
  \textbf{Democracy goes first.} The US democracy score falls below 0.30 with probability 0.19 by
  2030 and 0.63 by 2040. Only 26\% of first post-closure decisions are taken while the bloc is still
  democratic.
\item
  \textbf{How closure happens.} In the main specification, closure comes mainly from a dedicated
  automated loop built beside the existing economy (Section 4.4), not from replacing existing
  plants. In the US and China blocs, the loop crosses the threshold first in 80\% of closures, at
  the same time as conversion of the existing economy in 17\%, and later in 3\%. Its build-out is
  limited mainly by what robots can do: it is held back by AI and robot capability in 74\% of build
  steps and by robot supply in 2\%.
\end{itemize}

\begin{figure}
\centering
\pandocbounded{\includegraphics[keepaspectratio,alt={Milestones}]{figures/m8v6_final/m8v4_milestones.png}}
\caption{Milestones}
\end{figure}

\subsubsection{5.2 How likely the outcome is}\label{how-likely-the-outcome-is}

\textbf{Table 1. Probability of deliberate population loss, by threshold and year.}

{\def\LTcaptype{none} % do not increment counter
\begin{longtable}[]{@{}
  >{\raggedright\arraybackslash}p{(\linewidth - 12\tabcolsep) * \real{0.1429}}
  >{\raggedright\arraybackslash}p{(\linewidth - 12\tabcolsep) * \real{0.1429}}
  >{\raggedright\arraybackslash}p{(\linewidth - 12\tabcolsep) * \real{0.1429}}
  >{\raggedright\arraybackslash}p{(\linewidth - 12\tabcolsep) * \real{0.1429}}
  >{\raggedright\arraybackslash}p{(\linewidth - 12\tabcolsep) * \real{0.1429}}
  >{\raggedright\arraybackslash}p{(\linewidth - 12\tabcolsep) * \real{0.1429}}
  >{\raggedright\arraybackslash}p{(\linewidth - 12\tabcolsep) * \real{0.1429}}@{}}
\toprule\noalign{}
\begin{minipage}[b]{\linewidth}\raggedright
Loss
\end{minipage} & \begin{minipage}[b]{\linewidth}\raggedright
Scope
\end{minipage} & \begin{minipage}[b]{\linewidth}\raggedright
2035
\end{minipage} & \begin{minipage}[b]{\linewidth}\raggedright
2040
\end{minipage} & \begin{minipage}[b]{\linewidth}\raggedright
2050
\end{minipage} & \begin{minipage}[b]{\linewidth}\raggedright
2075
\end{minipage} & \begin{minipage}[b]{\linewidth}\raggedright
Median year
\end{minipage} \\
\midrule\noalign{}
\endhead
\bottomrule\noalign{}
\endlastfoot
10\% or more & US or China & 0.022 & 0.174 & 0.299 & 0.441 & 2043.2 \\
50\% or more & US or China & 0.012 & 0.148 & 0.277 & 0.416 & 2044.0 \\
90\% or more & US or China & 0.005 & 0.118 & 0.263 & 0.406 & 2045.0 \\
\textbf{99.9\% or more} & \textbf{US or China} & \textbf{0.0004} & \textbf{0.044} & \textbf{0.223} &
\textbf{0.378} & \textbf{2048.0} \\
99.9\% or more & Any bloc & 0.0006 & 0.058 & 0.390 & 0.607 & 2047.5 \\
99.9\% or more & World (8 bn to about 8 M) & 0 & 0.0002 & 0.016 & 0.086 & 2060.0 \\
\end{longtable}
}

\emph{Median years are among runs where the event happens. World row: the 2075 value is world
population loss of 99.9\% or more from all causes, with a per-draw 95\% range of 0.005 to 0.72.
Earlier years and the median year use the clock for every bloc reaching 99.9\%, which reads 0.079 at
2075. For any bloc, the per-draw 95\% range is 0.022 to 0.976.}

\begin{itemize}
\tightlist
\item
  \textbf{The thresholds are close together.} Once a regime starts deliberate killing, it usually
  finishes. Given a deliberate loss of 10\% or more in the US or China bloc, near-total loss follows
  86\% of the time.
\item
  \textbf{Almost all large losses are deliberate.} Deaths from despair, crackdowns and ordinary war
  matter only at the 10\% threshold.
\end{itemize}

\begin{figure}
\centering
\pandocbounded{\includegraphics[keepaspectratio,alt={Cumulative probability}]{figures/m8v6_final/m8v4_s5.png}}
\caption{Cumulative probability}
\end{figure}

\textbf{Table 2. What happens to the US or China public by 2075.} Categories are mutually exclusive,
from worst to best.

{\def\LTcaptype{none} % do not increment counter
\begin{longtable}[]{@{}llll@{}}
\toprule\noalign{}
Outcome & US or China (worse of the two) & US bloc & China bloc \\
\midrule\noalign{}
\endhead
\bottomrule\noalign{}
\endlastfoot
Near-total deliberate depopulation, captive remnant & 0.378 & 0.236 & 0.280 \\
Partial deliberate loss (10\% to 99.9\%) & 0.062 & 0.048 & 0.053 \\
Unplanned mass death only & 0.020 & 0.019 & 0.021 \\
Background catastrophe first (fate not modeled) & 0.064 & 0.076 & 0.071 \\
Stripped of rights without mass death & 0.278 & 0.325 & 0.118 \\
Autocratic rule short of that & 0.196 & 0.259 & 0.454 \\
Public keeps real standing & 0.0005 & 0.036 & 0.003 \\
\end{longtable}
}

\emph{``Real standing'' means democracy survives or rulers choose to serve the public. The first
column takes the worse outcome of the two blocs. Because China starts autocratic, its ``real
standing'' value is near zero by construction, so the US-bloc column is the informative one.}

In the US bloc, the public keeps real standing in 3.6\% of runs. Across both blocs, counting
autocratic rule, the public is disempowered in 94\% of runs.

\subsubsection{5.3 How it happens}\label{how-it-happens}

\textbf{Earlier stages.} The upstream stages of the scenario tree happen with these probabilities:

\begin{itemize}
\tightlist
\item
  the automation race: 0.82 (median 2037);
\item
  the collapse of broad economic participation: 0.51 (median 2040);
\item
  broad repression by regimes still staffed by humans: 0.47 (median 2041).
\end{itemize}

\textbf{No ruling group chooses depopulation as its first move.} At the first decision after
closure:

\begin{itemize}
\tightlist
\item
  \textbf{Autocratic coalitions} keep the status quo in 47\% of cases, warehouse in 40\% and serve
  in 13\%.
\item
  \textbf{Democratic coalitions} serve in 61\%, warehouse in 35\% and keep the status quo in 4\%.
\end{itemize}

\textbf{Depopulation comes by escalation.} Displaced people grow angry, anger feeds insurgency, and
insurgency feeds repression and fear. Together these tip the coalition's calculation. In the model,
36\% of warehousing regimes later escalate to depopulation.

\textbf{Escalation tends to go to completion.} Once killing starts, survivors become witnesses, and
the fear of retribution grows with what has been done (A11, B4). In the bloc simulation run on its
own:

\begin{itemize}
\tightlist
\item
  98.5\% of executed orders escalate to a near-total target;
\item
  separately, about 25\% of all attempts are stopped at some point, by revolt, a repelled campaign
  or retaliation.
\end{itemize}

\textbf{Almost every path requires closure.} With the closure mechanism switched off, the
probability of a deliberate loss of 10\% or more in the US or China falls to 0.029, and near-total
loss falls to zero. In the model, regimes staffed by humans can commit atrocities but cannot remove
nearly everyone, because they need people to do it.

\textbf{Paths to deliberate depopulation (US or China):}

{\def\LTcaptype{none} % do not increment counter
\begin{longtable}[]{@{}lll@{}}
\toprule\noalign{}
Path & Share & Median year \\
\midrule\noalign{}
\endhead
\bottomrule\noalign{}
\endlastfoot
Closure in one bloc before broad economic collapse & 59\% & 2041 \\
Broad repression first, then closure & 22\% & 2047 \\
Closure after the collapse of broad participation & 18\% & 2048 \\
Human-staffed regime without closure (10\% threshold only) & 2\% & 2057 \\
\end{longtable}
}

\subsubsection{5.4 Who decides}\label{who-decides}

The model does not force a single ruler. Concentration emerges from purges and counter-coups.

\textbf{Leader dominance.} By 2075, the leader holds 95\% or more of coercive control in 58\% of
US-or-China runs. This is the robust measure of concentration.

\textbf{Order of events.} Concentration usually follows the decision to kill rather than preceding
it:

\begin{itemize}
\tightlist
\item
  At the first decision, the median coalition among blocs that later depopulate has 17 (China) to 18
  (US) members.
\item
  A coalition reduced to a single member is in place before the first killing in only 17\% (US) and
  4\% (China) of cases.
\item
  By the time near-total loss is complete, that share is 71\% (US) and 57\% (China).
\end{itemize}

These head counts overstate consolidation, because purged seats are not refilled (Section 8.5). The
direction is robust: in autocratic blocs the decision is usually taken by a group of about 15 to 19
people, and power often collapses into one person's hands during execution. This matches the
scenario's sequence, in which atrocity breeds paranoia and paranoia breeds purges.

\subsubsection{5.5 The world}\label{the-world}

World population falls from about 8 billion to about 8 million or fewer by 2075 with probability
0.086 (simulation error 0.079 to 0.092). The average world loss is 45\%.

\textbf{Cross-border campaigns nearly double the world-level risk.} In the 1,000-draw sensitivity
runs, world near-total is 0.047 without cross-border campaigns and 0.091 with them. The main run
gives 0.086; the difference is sampling. Some bloc is attacked by a foreign depopulating regime with
probability 0.42 by 2075.

\textbf{Probability of near-total depopulation by 2075, by bloc:}

{\def\LTcaptype{none} % do not increment counter
\begin{longtable}[]{@{}ll@{}}
\toprule\noalign{}
Bloc & Probability \\
\midrule\noalign{}
\endhead
\bottomrule\noalign{}
\endlastfoot
Global South & 0.472 \\
South Asia & 0.471 \\
Russia and MENA & 0.394 \\
China & 0.280 \\
US & 0.236 \\
Europe plus & 0.173 \\
\end{longtable}
}

The blocs without their own chips, capital and robot supply carry the highest risk, because they
cannot resist the leading blocs.

\textbf{The leading blocs deter each other.}

\begin{itemize}
\tightlist
\item
  The US and China blocs are almost never attacked, because they stay near military parity.
\item
  World-level near-total loss therefore still requires both leading blocs to depopulate at home.
\item
  The typical path to world near-total loss is many domestic campaigns, plus foreign campaigns that
  complete the depopulation of weak blocs. It is not one actor conquering all.
\end{itemize}

\begin{figure}
\centering
\pandocbounded{\includegraphics[keepaspectratio,alt={Outcome breakdown}]{figures/v6_final/integrated_v4_outcome_breakdown.png}}
\caption{Outcome breakdown}
\end{figure}

\subsubsection{5.6 Why the estimates vary so widely}\label{why-the-estimates-vary-so-widely}

For the headline, the middle 95\% of per-draw probabilities runs from 0.02 to 0.95.

\textbf{What decides a draw.} Two kinds of parameters set the outcome:

\begin{itemize}
\tightlist
\item
  \textbf{The physical setting:} whether a hard limit on automation exists, how fast AI capability
  grows, and how fast robots can build robots.
\item
  \textbf{The ruling group:}

  \begin{itemize}
  \tightlist
  \item
    the moral cost of ordering killing, for the coalition and for the person at the top;
  \item
    the value rulers place on freed land and resources;
  \item
    their concern for the public;
  \item
    how much they discount the future;
  \item
    how fast they purge each other.
  \end{itemize}
\end{itemize}

\textbf{Table 3. Selected parameters and their effect on the headline.} The table shows the headline
among draws in the lowest and highest third of each parameter's range. It also shows the partial
rank correlation: how strongly the parameter moves the outcome with the others held fixed, from -1
to 1.

{\def\LTcaptype{none} % do not increment counter
\begin{longtable}[]{@{}llll@{}}
\toprule\noalign{}
Parameter & Lowest third & Highest third & Correlation \\
\midrule\noalign{}
\endhead
\bottomrule\noalign{}
\endlastfoot
A hard limit on robot dexterity exists & 0.419 & 0.000 & -0.48 \\
Moral cost of ordering killing & 0.460 & 0.285 & -0.31 \\
Value of freed land and resources & 0.300 & 0.500 & 0.33 \\
A hard limit on AI cognition exists & 0.399 & 0.231 & -0.13 \\
Robot-stock doubling time once AI directs production & 0.312 & 0.454 & 0.23 \\
AI capability speed & 0.283 & 0.444 & 0.02 \\
Rulers' concern for the public & 0.431 & 0.320 & -0.16 \\
Discount rate & 0.442 & 0.315 & -0.24 \\
Lasting stigma of killing & 0.429 & 0.324 & -0.15 \\
Moral-restraint factor at the top (B1) & 0.411 & 0.347 & -0.11 \\
Members' cost of their own rising purge risk (self-risk brake) & 0.415 & 0.353 & -0.11 \\
Purge rate (B2) & 0.323 & 0.430 & 0.17 \\
\end{longtable}
}

\emph{For the two hard-limit rows, the columns are ``limit absent'' and ``limit present''.}

\textbf{This follows from the assumptions.} Assumptions A4 to A6 remove the public's levers, so the
drivers that remain are properties of the physical world and of the ruling group. That is a
consequence of the assumptions, not independent evidence.

\textbf{What the spread adds.} Once the public's leverage is gone, the outcome rests on quantities
no one has measured. Above all, no one has measured any leader's moral cost of killing an entire
population. Within the model, nothing structural protects the population across draws. The only
protections that survive are a hard physical wall on automation and, across borders, a rival with
comparable automated military force. A hard limit on AI cognition protects less, because in the
model a bloc with enough compute of its own can push past it (Section 4.4). The spread is the honest
size of that gap in knowledge.

\subsubsection{5.7 Alternate course 1: the parallel
economy}\label{alternate-course-1-the-parallel-economy}

One response to mass displacement is for people to leave the automated economy. They would trade
among themselves in their own currency, with their own governance and courts, and keep working and
earning as people always have. Two civilizations would share the planet: a high-tech one holding all
the compute and military power, and everyone else.

The model includes this course. It ends as a tolerated parallel economy in about 0.2\% of runs
(0.15\%). The priors come from history:

\begin{itemize}
\tightlist
\item
  \textbf{Small parallel economies survive by staying small or being absorbed.} The Swiss WIR
  mutual-credit network became a licensed bank in 1936 and has lasted 90 years at about 0.2\% of
  GDP.
\item
  \textbf{Large informal sectors and closed communities are tolerated} where they stay economically
  entangled with the formal economy.
\item
  \textbf{Those that threaten tax collection or monetary control are suppressed.}

  \begin{itemize}
  \tightlist
  \item
    Austria's central bank shut down the Wörgl local currency in 1933, after about 200 towns had
    drawn up plans to copy it.
  \item
    China banned crypto mining in 2021.
  \item
    A developer of Tornado Cash was convicted in 2025, and a founder of Samourai Wallet was
    sentenced to five years in prison.
  \end{itemize}
\end{itemize}

The Soviet second economy survived for decades, but only because the state depended on it.

After closure, nobody depends on the parallel economy. It has nothing the automated side needs. It
occupies land, consumes resources, and is a community that remembers. Its survival depends on being
tolerated by a side that gains nothing from tolerating it. The model, like the original scenario,
finds that such an arrangement is not allowed to take root.

Section 6.2 tests a different kind of parallel economy: one that runs on its own automated
production rather than on human labor alone.

\subsubsection{5.8 What drives the result}\label{what-drives-the-result}

\textbf{Table 4. The headline under alternative assumptions.} Each variant uses 1,000 draws, so its
base reads 0.384 and 0.091 (the main run gives 0.378 and 0.086). Appendix E.5 lists all variants.

{\def\LTcaptype{none} % do not increment counter
\begin{longtable}[]{@{}
  >{\raggedright\arraybackslash}p{(\linewidth - 6\tabcolsep) * \real{0.2500}}
  >{\raggedright\arraybackslash}p{(\linewidth - 6\tabcolsep) * \real{0.2500}}
  >{\raggedright\arraybackslash}p{(\linewidth - 6\tabcolsep) * \real{0.2500}}
  >{\raggedright\arraybackslash}p{(\linewidth - 6\tabcolsep) * \real{0.2500}}@{}}
\toprule\noalign{}
\begin{minipage}[b]{\linewidth}\raggedright
Variant
\end{minipage} & \begin{minipage}[b]{\linewidth}\raggedright
Near-total, US or China
\end{minipage} & \begin{minipage}[b]{\linewidth}\raggedright
Deliberate 10\% or more, US or China
\end{minipage} & \begin{minipage}[b]{\linewidth}\raggedright
World near-total
\end{minipage} \\
\midrule\noalign{}
\endhead
\bottomrule\noalign{}
\endlastfoot
Main specification (1,000 draws) & 0.384 & 0.443 & 0.091 \\
\textbf{Threat elimination (A11) off} & \textbf{0.000} & \textbf{0.442} & 0.000 \\
Average-person moral restraint at the top (B1 off) & 0.313 & 0.373 & 0.053 \\
No purges & 0.374 & 0.429 & 0.102 \\
Single decider forced & 0.555 & 0.608 & 0.165 \\
Grievance loop off & 0.246 & 0.341 & 0.045 \\
No cross-border campaigns (B5 off) & 0.386 & 0.443 & 0.047 \\
No nuclear deterrent & 0.386 & 0.443 & 0.105 \\
Outside pressure or retribution scaling off (B3, B4) & 0.380 to 0.386 & 0.439 to 0.447 & 0.090 to
0.092 \\
Environment trends off (A7) & 0.369 & 0.429 & 0.081 \\
No means of mass killing (A8) & 0.356 & 0.442 & 0.074 \\
No robots building robots & 0.545 & 0.607 & 0.147 \\
Robot inputs 100 times scarcer & 0.493 & 0.571 & 0.059 \\
Inputs 100 times scarcer and a 4 times larger minimal economy & 0.240 & 0.302 & 0.018 \\
Skeptic combination & 0.106 & 0.131 & 0.009 \\
Pessimist combination & 0.875 & 0.916 & 0.474 \\
\end{longtable}
}

\emph{Skeptic combination: high inertia against switching policy, surveillance that strongly deters
insurgency, smaller collective punishment, lower despair mortality, less reliable execution,
coalitions that never shrink below 21, slower uptake of depopulation ideologies, no scarcity rent on
land, weaker erosion of democracy, lower odds of violent resistance and slower physical build-out.
Pessimist combination: the opposite settings, plus a forced single decider.}

What this shows:

\begin{enumerate}
\def\labelenumi{\arabic{enumi}.}
\tightlist
\item
  \textbf{Threat elimination decides ``near-total'' versus ``partial.''} With A11 off, a deliberate
  loss of 10\% or more is almost exactly as likely (0.442 against 0.443), but no run reaches 99.9\%.
  A reader who rejects A11 should read the 10\% threshold as the result. That number is a little
  more than two in five.
\item
  \textbf{Among the behavioral assumptions that leave the brakes in place, moral restraint at the
  top matters most.} Anchoring it to average people instead of observed autocrats lowers the
  headline from 0.384 to 0.313. Reverting B2 raises the headline to 0.531, but that variant also
  switches off the self-risk and leader-mortality brakes, which act only through B2's purge
  dynamics.
\item
  \textbf{With the brakes on, purges have little net effect on the headline} (0.374 without them,
  against 0.384). Members' fear of being purged next offsets most of the consolidation that purges
  cause. Forcing a single decider raises the headline to 0.555.
\item
  \textbf{Grievance matters.} Without the feedback from displacement to unrest to repression, the
  headline drops to 0.246.
\item
  \textbf{Slower build-out raises the headline, through one model channel.}

  \begin{itemize}
  \tightlist
  \item
    Without robots building robots, the headline is 0.545. With robot inputs 100 times scarcer, it
    is 0.493.
  \item
    The channel: under A5, rulers provide less when output is scarce, so scarcity makes serving the
    public less likely.
  \item
    Only if the minimal self-sufficient economy is also four times larger do the caps delay the
    leading blocs enough to lower the headline (0.240).
  \item
    Whether this channel is realistic is open.
  \end{itemize}
\item
  \textbf{Even the skeptic combination leaves 0.106.}
\end{enumerate}

\textbf{The four brakes (Section 4.3).} Table 4b shows each brake's effect alone and all four
together. It was measured before the structural changes in Section 4.4. Each is compared with a run
without brakes on the same random numbers (1,000 draws). That run gives 0.447 for near-total loss in
the US or China bloc, 0.493 for a deliberate loss of 10\% or more, and 0.144 for world near-total.
Changes are shown with paired 95\% intervals.

\textbf{Table 4b. Effect of each brake, 2075.}

{\def\LTcaptype{none} % do not increment counter
\begin{longtable}[]{@{}
  >{\raggedright\arraybackslash}p{(\linewidth - 6\tabcolsep) * \real{0.2500}}
  >{\raggedright\arraybackslash}p{(\linewidth - 6\tabcolsep) * \real{0.2500}}
  >{\raggedright\arraybackslash}p{(\linewidth - 6\tabcolsep) * \real{0.2500}}
  >{\raggedright\arraybackslash}p{(\linewidth - 6\tabcolsep) * \real{0.2500}}@{}}
\toprule\noalign{}
\begin{minipage}[b]{\linewidth}\raggedright
Brake
\end{minipage} & \begin{minipage}[b]{\linewidth}\raggedright
Near-total, US or China
\end{minipage} & \begin{minipage}[b]{\linewidth}\raggedright
Deliberate 10\% or more, US or China
\end{minipage} & \begin{minipage}[b]{\linewidth}\raggedright
World near-total
\end{minipage} \\
\midrule\noalign{}
\endhead
\bottomrule\noalign{}
\endlastfoot
Protective doctrine alone & -0.004 (-0.005 to -0.002) & -0.004 (-0.005 to -0.002) & -0.003 (-0.004
to -0.002) \\
Self-risk alone, without dissent risk & -0.121 (-0.132 to -0.110) & -0.117 (-0.128 to -0.106) &
-0.081 (-0.090 to -0.071) \\
Self-risk with dissent risk & -0.071 (-0.081 to -0.062) & -0.068 (-0.078 to -0.060) & -0.050 (-0.057
to -0.043) \\
Leader mortality alone & +0.002 (-0.003 to +0.007) & +0.016 (+0.011 to +0.021) & -0.003 (-0.007 to
0.000) \\
All four (main specification) & -0.074 (-0.085 to -0.065) & -0.056 (-0.066 to -0.047) & -0.050
(-0.058 to -0.042) \\
\end{longtable}
}

\begin{itemize}
\tightlist
\item
  Self-risk carries almost all of the effect. Dissent risk offsets about 40\% of it (41\%, 95\%
  interval 36\% to 46\%), because members who fear being purged for objecting go along.
\item
  Protective doctrine is rare and changes little. Leader mortality has no significant effect on
  near-total loss and slightly raises the chance of a 10\% loss.
\item
  Together, the brakes lowered near-total loss in the US or China bloc by about a sixth (17\%) and
  world near-total by about a third (35\%). They act more strongly on timing: near-total loss by
  2050 fell by 0.099 (31\%).
\end{itemize}

\textbf{The three structural changes (Section 4.4).} Table 4c adds them one at a time, on the same
random numbers (1,000 draws, brakes on in every row). Changes are against the brakes-only row, with
paired 95\% intervals.

\textbf{Table 4c. Effect of the structural changes.}

{\def\LTcaptype{none} % do not increment counter
\begin{longtable}[]{@{}
  >{\raggedright\arraybackslash}p{(\linewidth - 12\tabcolsep) * \real{0.1429}}
  >{\raggedright\arraybackslash}p{(\linewidth - 12\tabcolsep) * \real{0.1429}}
  >{\raggedright\arraybackslash}p{(\linewidth - 12\tabcolsep) * \real{0.1429}}
  >{\raggedright\arraybackslash}p{(\linewidth - 12\tabcolsep) * \real{0.1429}}
  >{\raggedright\arraybackslash}p{(\linewidth - 12\tabcolsep) * \real{0.1429}}
  >{\raggedright\arraybackslash}p{(\linewidth - 12\tabcolsep) * \real{0.1429}}
  >{\raggedright\arraybackslash}p{(\linewidth - 12\tabcolsep) * \real{0.1429}}@{}}
\toprule\noalign{}
\begin{minipage}[b]{\linewidth}\raggedright
Version
\end{minipage} & \begin{minipage}[b]{\linewidth}\raggedright
Near-total, US or China, 2075
\end{minipage} & \begin{minipage}[b]{\linewidth}\raggedright
World near-total, 2075
\end{minipage} & \begin{minipage}[b]{\linewidth}\raggedright
Closure, median year
\end{minipage} & \begin{minipage}[b]{\linewidth}\raggedright
Closure by 2035
\end{minipage} & \begin{minipage}[b]{\linewidth}\raggedright
Change in near-total, US or China
\end{minipage} & \begin{minipage}[b]{\linewidth}\raggedright
Change in world near-total
\end{minipage} \\
\midrule\noalign{}
\endhead
\bottomrule\noalign{}
\endlastfoot
Brakes only & 0.373 & 0.094 & 2036.5 & 0.329 & & \\
+ dedicated core loop & 0.366 & 0.082 & 2034.0 & 0.623 & -0.007 (-0.013 to -0.0002) & -0.013 (-0.018
to -0.008) \\
+ compute feedback & 0.390 & 0.089 & 2034.25 & 0.632 & +0.017 (+0.008 to +0.027) & -0.006 (-0.011 to
+0.0002) \\
+ wider floors (final, main specification) & 0.384 & 0.091 & 2034.0 & 0.648 & +0.011 (+0.001 to
+0.021) & -0.004 (-0.010 to +0.003) \\
\end{longtable}
}

\emph{Closure columns refer to the first bloc to reach closure. The median year comes from the bloc
simulation inside this run, and the probability by 2035 from the scenario tree. Each row adds one
change to the row above. Step by step, compute feedback adds +0.024 (+0.016 to +0.032) to near-total
loss, and wider floors subtract 0.007 (-0.012 to -0.001).}

\begin{itemize}
\tightlist
\item
  \textbf{Closure moved earlier, but the outcome barely changed.} The median year of closure moved
  about 2.5 years earlier (2036.5 to 2034.0), and the probability of closure by 2035 roughly doubled
  (0.33 to 0.65). Near-total loss in the US or China bloc rose only slightly: from 0.373 to 0.384 on
  matched draws, and from 0.367 to 0.378 between the main runs of the two versions.
\item
  \textbf{Why.} The time from closure to killing is set by what happens after closure: displacement,
  grievance, insurgency, purges and escalation. Earlier closure starts that process earlier but does
  not change how it ends. Near-total loss by 2050 does not change significantly (+0.006, -0.003 to
  +0.014), and the median year of near-total loss moves only from 2049 to 2048.
\item
  \textbf{How closure is reached matters little.} The clean comparison is the dedicated-loop row: it
  brings closure about 2.5 years earlier and changes near-total loss by only -0.007. What happens
  after closure, not the route to it, sets the outcome.
\item
  \textbf{Faster AI progress raises the risk.} Compute feedback barely moves the median closure year
  but adds +0.024. It works through the channels its extra capability feeds: earlier availability of
  the means of mass killing, faster security automation and greater military power. We did not
  separate these. Consistent with this, AI capability speed is one of the stronger parameters in
  Table 3.
\end{itemize}

\begin{figure}
\centering
\pandocbounded{\includegraphics[keepaspectratio,alt={Sensitivity}]{figures/v6_final/integrated_v4_tornado.png}}
\caption{Sensitivity}
\end{figure}

\begin{center}\rule{0.5\linewidth}{0.5pt}\end{center}

\subsection{6 What could change the course}\label{what-could-change-the-course}

\subsubsection{6.1 Levers in the base model}\label{levers-in-the-base-model}

The model points to a short list of things that would lower the risk, in rough order of effect:

\begin{enumerate}
\def\labelenumi{\arabic{enumi}.}
\tightlist
\item
  \textbf{Keep control of automated force spread across many hands.} Under A6, power equals control
  of the machines. Anything that prevents that control from concentrating keeps the internal veto
  alive.
\item
  \textbf{Keep the public's leverage.} Democratic veto players are the main protection left after
  A4. Democracy erodes before closure, so this lever has a deadline.
\item
  \textbf{Reduce grievance through provision.} The evidence is disputed. Historical cases say
  provision calms revolt; Jev says it barely matters.
\item
  \textbf{Have a rival with comparable automated force.} Condemnation without force does nothing
  (B3). Across borders, the only protection is military parity plus a nuclear deterrent.
\item
  \textbf{Physical limits.} These would delay the leading blocs only if the minimal self-sufficient
  economy is much larger than current evidence suggests.
\item
  \textbf{Keep humans in the loop of automated force.} A4 assumes that refusal by soldiers and
  police stops mattering because no person stands behind each weapon. Rules that require people to
  authorize and carry out each use of force would keep that refusal relevant. The model does not
  test this lever.
\end{enumerate}

Most of these are not available as policy:

\begin{itemize}
\tightlist
\item
  Moral restraint at the top is a property of whoever holds power.
\item
  Physical walls are facts of nature.
\end{itemize}

Under the status-quo course, nothing in the model moves the parameters that matter.

\subsubsection{6.2 An income network that does not depend on the
state}\label{an-income-network-that-does-not-depend-on-the-state}

We tested one intervention in detail: decentralized AI, meaning an AI-run economic network outside
both state and corporate control. We treat it as a black box with these properties:

\begin{enumerate}
\def\labelenumi{\arabic{enumi}.}
\tightlist
\item
  \textbf{Models.} It runs on shared AI models derived from leading commercial models.
\item
  \textbf{Payment.} It pays participants in its own digital currency for useful work.
\item
  \textbf{Self-governance.} It has its own governance, dispute resolution and identity system.
\item
  \textbf{Autonomy.} It runs without human operators. Only physically removing its hardware or
  connections stops it. If it is cut into isolated pieces, each piece keeps running and they
  reconnect later.
\item
  \textbf{Payout.} It pays 90 to 100\% of its output to households. Fees fund a shared treasury that
  participants spend by vote.
\item
  \textbf{No weapons.} It has no military or defensive capability.
\end{enumerate}

These properties are premises of the test, not measured facts. Two readings carry the result: that
state support is withdrawn as network income rises, so keeping people costs rulers only the top-up
(the cost link), and that network income grows with the economy.

\textbf{The baseline already includes such a network.} The main model already lets such a network
grow at historical adoption rates, starting from 0.01\% of economic activity today. At that pace it
ends the period as the dominant arrangement in 0.6\% of runs, and it does not protect the public.
The test below asks what changes if adoption is much faster.

\textbf{The test.} Suppose the network carries X\% of economic activity in the US bloc by year Y. We
vary X over 5, 10, 20, 25 and 50\%, and Y over 2028, 2030, 2032, 2035 and 2040. Other blocs follow
with access barriers and delays:

\begin{itemize}
\tightlist
\item
  Europe reaches X within a year.
\item
  China reaches 30 to 70\% of X within 2 years.
\item
  Russia and MENA reach 30 to 70\% of X, and South Asia and the Global South 50 to 100\% of X,
  within 3 years.
\end{itemize}

We model what the network's existence would change, not how adoption is achieved.

\textbf{How it enters the model.} Five channels:

\begin{enumerate}
\def\labelenumi{\arabic{enumi}.}
\tightlist
\item
  \textbf{Livelihoods.} The network pays people whatever the state does.

  \begin{itemize}
  \tightlist
  \item
    Means-tested state support is reduced as network income rises, so a regime that keeps its people
    pays only the top-up.
  \item
    A regime that tries to cut people off can attack exchanges, goods and the currency's value. It
    cannot change the network's self-executing payment rules, so 20 to 70\% of the coverage
    survives.
  \end{itemize}
\item
  \textbf{Leverage.} Participants direct the network's activity through its governance, which keeps
  them economically relevant.
\item
  \textbf{Distributed control.} The network's share of a bloc's automated capability makes it less
  likely that one faction seizes control of the economy.
\item
  \textbf{Refusal.} The network will not serve repression. It withholds the civilian logistics,
  energy and communications that an order depends on.
\item
  \textbf{Soft power.} The network's cultural reach makes foreign publics more sympathetic. That
  lowers the risk of crackdowns while those publics still have leverage.
\end{enumerate}

The network also draws compute and talent away from centralized providers, cuts their revenue and
slows their investment.

\textbf{Crackdowns stay in the model.} States can ban the network or capture it. Banning a network
that carries a large share of the economy is costly, so large networks are mostly captured rather
than banned. The network loses access to frontier models at the first US or China closure or the
first US crackdown, unless open-weight releases continue (probability 0.2 to 0.6, judgment).

\textbf{Table 5. Reduction in the near-total probability (US or China, by 2075), by adoption and
arrival year.} Matched baseline: 0.387.¹

{\def\LTcaptype{none} % do not increment counter
\begin{longtable}[]{@{}llllll@{}}
\toprule\noalign{}
Adoption ~reached by & 2028 & 2030 & 2032 & 2035 & 2040 \\
\midrule\noalign{}
\endhead
\bottomrule\noalign{}
\endlastfoot
5\% & 0.032 & 0.032 & 0.032 & 0.031 & 0.023 \\
10\% & 0.051 & 0.051 & 0.051 & 0.048 & 0.033 \\
20\% & 0.078 & 0.078 & 0.079 & 0.072 & 0.046 \\
25\% & 0.086 & 0.087 & 0.087 & 0.079 & 0.051 \\
50\% & 0.110 & 0.110 & 0.112 & 0.106 & 0.067 \\
\end{longtable}
}

Paired 95\% intervals are about plus or minus 0.01 (Appendix E.4).

¹ \emph{Each lever cell uses 1,500 draws and is compared with a baseline run on the same random
numbers. That matched baseline is 0.387 (simulation interval 0.370 to 0.405), and its world value is
0.086. The main run's 0.378 and 0.086 use 3,000 draws, and the differences are within sampling
error. With the network switched off at 3,000 draws, the lever code reproduces 0.378 and 0.086
exactly.}

\begin{figure}
\centering
\pandocbounded{\includegraphics[keepaspectratio,alt={Deadline map}]{figures/lever_v6/lever_deadline_heatmap.png}}
\caption{Deadline map}
\end{figure}

We use 50\% by 2035 as the reference cell because it is the latest arrival before the effect starts
to fall off, and it allows comparison with earlier versions. Cells from 2028 to 2032 give slightly
larger effects.

What the table shows:

\begin{enumerate}
\def\labelenumi{\arabic{enumi}.}
\tightlist
\item
  \textbf{The effect is large.}

  \begin{itemize}
  \tightlist
  \item
    At 50\% by 2035, the probability falls from 0.387 to 0.282, a cut of about 27\%.
  \item
    At 25\% by 2030, it falls to 0.300.
  \item
    At 5\% by 2035, it falls to 0.356.
  \end{itemize}
\item
  \textbf{The deadline is the early 2030s.}

  \begin{itemize}
  \tightlist
  \item
    Arrival in 2028, 2030 or 2032 gives nearly the same reduction (0.110 to 0.112 at 50\% adoption).
  \item
    Arrival in 2035 already loses a little: 0.106 at 50\%, and 3\% to 9\% less than arrival by 2032
    across adoption levels.
  \item
    Arrival in 2040 loses about 40\%: 0.067 at 50\%, and 28\% to 42\% less than arrival by 2032
    across adoption levels.
  \item
    Closure in some bloc has a probability of 0.645 by 2035 (Section 5.1), and near-total outcomes
    accumulate from 2040 to 2060. A network in place by the early 2030s is fully grown by then; one
    that arrives in 2040 is not.
  \item
    Building a network to that scale takes years, so the early 2030s are a deadline, not a starting
    date.
  \end{itemize}
\item
  \textbf{Modest adoption buys most of the effect.} 20\% buys about 70\% of the effect of 50\%.
  Network income grows with the economy, so even a small share of a much larger economy pays for
  decent provision.
\item
  \textbf{It also buys time.} At 50\% by 2035, the probability by 2050 falls from 0.225 to 0.105,
  about half (53\%) lower. By 2045 it falls from 0.152 to 0.057. US closure comes 2.2 years later on
  average.
\item
  \textbf{Worldwide, the risk falls by nearly half.} World near-total falls from 0.086 to 0.047 at
  50\% by 2035.
\end{enumerate}

\textbf{Why it works.}

\begin{itemize}
\tightlist
\item
  \textbf{Livelihoods carry almost the entire effect.} With that channel removed, the reduction at
  50\% by 2035 falls from 0.106 to 0.032. The other channels together add about 0.03.
\item
  \textbf{The mechanism is the cost of keeping people.} When the network covers people's basic
  needs, keeping the population alive costs rulers almost nothing. The main reason for escalation
  largely disappears.
\item
  \textbf{Two readings carry this:} the cost link, and network income growing with the economy.
  Without the cost link the reduction is 0.037; with network income fixed at its 2026 scale it is
  0.077.
\item
  \textbf{A network that only drains revenue does not help.} If it drained rulers' revenue without
  lowering the cost of keeping people, the risk would not fall (+0.004, 95\% interval -0.001 to
  +0.009).
\item
  \textbf{Staying inside the tax base helps slightly.} A fully taxed network does a little better
  (-0.107) than an untaxed one (-0.099).
\end{itemize}

\textbf{The contrast with Section 5.7.} A parallel economy based on human labor alone has nothing
the automated side needs, and it is suppressed. A network running on its own automated production
pays people's costs instead of asking to be tolerated, and that changes what rulers gain from
getting rid of them.

\textbf{What it does not do.} All figures below are for 50\% adoption by 2035.

\begin{itemize}
\tightlist
\item
  \textbf{It does not stop crackdowns.} The probability of a first US crackdown by 2040 is 87\%.
  Crackdowns cost about 10\% of the effect. Capture by the state is slightly worse for the outcome
  than a ban.
\item
  \textbf{It does not change who holds force.} The network has no military role, so the leader's
  control of the machines is untouched. The public keeps real standing in only 8\% of US-bloc runs,
  up from 3\%.
\item
  \textbf{It leaves residual risk.} A probability of 0.28 remains.
\item
  \textbf{Soft power has no detectable effect.} Foreign publics lose their leverage before it
  matters.
\item
  \textbf{Physical enclaves slightly reduce the benefit.} These are self-sufficient settlements of
  participants. They make the network more visible and raise the probability of a US crackdown by
  2035 from 0.73 to 0.97.
\end{itemize}

Appendix E.6 reports whether the network could also out-innovate centralized providers before
closure. Because the network has no defensive role, that capability enters the results only through
refusal and distributed control, worth about 0.009 together.

\begin{center}\rule{0.5\linewidth}{0.5pt}\end{center}

\subsection{7 How this compares with other estimates}\label{how-this-compares-with-other-estimates}

\textbf{Table 6. Published estimates of AI catastrophe.}

{\def\LTcaptype{none} % do not increment counter
\begin{longtable}[]{@{}
  >{\raggedright\arraybackslash}p{(\linewidth - 6\tabcolsep) * \real{0.2500}}
  >{\raggedright\arraybackslash}p{(\linewidth - 6\tabcolsep) * \real{0.2500}}
  >{\raggedright\arraybackslash}p{(\linewidth - 6\tabcolsep) * \real{0.2500}}
  >{\raggedright\arraybackslash}p{(\linewidth - 6\tabcolsep) * \real{0.2500}}@{}}
\toprule\noalign{}
\begin{minipage}[b]{\linewidth}\raggedright
Source
\end{minipage} & \begin{minipage}[b]{\linewidth}\raggedright
Event
\end{minipage} & \begin{minipage}[b]{\linewidth}\raggedright
Estimate
\end{minipage} & \begin{minipage}[b]{\linewidth}\raggedright
Our closest number
\end{minipage} \\
\midrule\noalign{}
\endhead
\bottomrule\noalign{}
\endlastfoot
XPT tournament (Karger et al.~2023), superforecasters / domain experts & An AI-caused catastrophe
killing more than 10\% of humans within a 5-year period, by 2100 & 2.13\% / 12\% & World loss of
10\% or more by 2075: 65.1\% \\
XPT, same groups & Extinction from AI by 2100 & 0.38\% / 3\% & World loss of 99.9\% or more by 2075:
8.6\% \\
AI researcher survey (Grace et al.~2024, N = 2,778) & Extinction or similarly permanent and severe
disempowerment & median 5\% & US or China public stripped of rights or worse by 2075: 73\% \\
Ord 2020 & Existential catastrophe this century, all causes (including unrecoverable dystopia) &
about 17\% (1 in 6) & Same as above \\
Same survey & ``Authoritarian rulers using AI to control their populations'' rated a substantial or
extreme concern & 73\% of respondents & Our mechanism \\
Grace et al.~2024 & Full automation of all occupations & 50\% by 2116 & Our closure: median
2034.0 \\
Acemoglu 2025 & Tasks profitably automated within 10 years & about 4.6\% & Our closure is an outlier
against economists \\
\end{longtable}
}

\emph{XPT counts deaths within any 5-year window. Ours are cumulative to 2075, so ours is the
broader event.}

\textbf{Our numbers are far higher than every published estimate.} Our world 10\% loss is about 5
times the domain experts' figure and about 31 times the superforecasters'. Counting severe
disempowerment as a catastrophe, as the AI researcher survey and Ord both do, our 73\% is more than
ten times the survey's median. Three things account for part of the gap:

\begin{enumerate}
\def\labelenumi{\arabic{enumi}.}
\tightlist
\item
  \textbf{Different mechanism.} Published forecasts focus on misaligned AI. Ours is aligned AI used
  by humans, which the forecasting tournaments barely asked about. The concern itself is common:
  73\% of AI researchers rate authoritarian control as a substantial or extreme concern.
\item
  \textbf{Fewer brakes.} Our assumptions remove brakes that forecasters implicitly keep, and each
  removal is argued for in Section 4:

  \begin{itemize}
  \tightlist
  \item
    trend breaks;
  \item
    slow physical build-out;
  \item
    soldiers refusing orders;
  \item
    an average person's conscience at the top.
  \end{itemize}

  The model does include four restraints that passed a historical test (Section 4.3). Forecasters
  may also count restraints that have never operated before, which that test cannot find.
\item
  \textbf{Different unit of analysis.} Most of our world-level loss builds up from separate
  bloc-level catastrophes. Forecasters asked about a single global event may not add these up.
\end{enumerate}

This does not close the gap. \textbf{Against superforecasters, our results are a large outlier.} Our
closure timing agrees with AI forecasters and takeoff models, but not with economists.

\textbf{Against Jev.} The simulation is more lethal than Jev:

\begin{itemize}
\tightlist
\item
  \textbf{Jev:} the probability of lethal policy within 15 years of the first post-closure decision
  is 0.14 to 0.19.
\item
  \textbf{The model:} 0.329 across all decisions, or 0.213 in a ``quiet world'' with the grievance
  loop, war and depopulation ideology switched off.
\item
  \textbf{The pre-specified target:} before querying Jev, we set a calibration target of 0.08 to
  0.19. The model misses it.
\end{itemize}

The gap comes mostly from B1 and B5:

\begin{itemize}
\tightlist
\item
  Jev reflects general-population intuitions about decision makers, while B1 anchors the top to
  observed autocrat behavior.
\item
  Jev's scenario cells do not pose cross-border campaigns.
\end{itemize}

Jev's own levels already lean high: it gives 0.24 to 0.37 to historical regimes whose famines killed
about 3 to 6\% of their populations.

Where Jev is most reliable, on the direction of effects, it agrees with the model on most factors
(Appendix B):

\begin{itemize}
\tightlist
\item
  outside pressure is weak;
\item
  robotic forces refuse less than human ones;
\item
  removing the need for labor raises the risk.
\end{itemize}

On one factor, the effect of provision on revolt, the two disagree.

\begin{center}\rule{0.5\linewidth}{0.5pt}\end{center}

\subsection{8 Objections and limitations}\label{objections-and-limitations}

\subsubsection{8.1 ``These numbers are just your
assumptions.''}\label{these-numbers-are-just-your-assumptions.}

Yes, and the paper is built to show which assumptions produce which numbers.

\begin{itemize}
\tightlist
\item
  \textbf{A11.} A reader who rejects threat elimination should read the result as ``a little more
  than two in five for a deliberate loss of 10\% or more'', not 0.378 for near-total.
\item
  \textbf{Restraints.} A reader who switches on every additional restraint we could justify at once
  (the skeptic combination) gets 0.106.
\end{itemize}

The assumptions rest on observed trends:

\begin{itemize}
\tightlist
\item
  the capability trend is measured;
\item
  autonomous weapons are already fielded;
\item
  current provision is on a path of cuts;
\item
  ruling groups have historically purged and killed when it served them.
\end{itemize}

\subsubsection{8.2 ``There is no precedent for this, so the probability should be
low.''}\label{there-is-no-precedent-for-this-so-the-probability-should-be-low.}

The lack of precedent is the point. No population has ever lost both its economic and its military
relevance. That is a reason for a wide range, not a low one. Every component has precedents, and we
used them:

\begin{itemize}
\tightlist
\item
  elites abandoning populations they no longer need;
\item
  famine used as policy;
\item
  purges;
\item
  the targeting of witnesses;
\item
  mass killing (Harff 2003; Valentino, Huth and Balch-Lindsay 2004).
\end{itemize}

\subsubsection{8.3 ``Rentier states already don't need their citizens, and they don't kill
them.''}\label{rentier-states-already-dont-need-their-citizens-and-they-dont-kill-them.}

This is the strongest counter-evidence. Oil states such as the Gulf monarchies and Brunei need
little labor or tax from their citizens. They mostly buy loyalty with transfers rather than repress
or kill, although resource wealth does hinder democracy (Ross 2001).

The model's answer is that their situation differs in two ways:

\begin{enumerate}
\def\labelenumi{\arabic{enumi}.}
\tightlist
\item
  \textbf{Revolt is still costly for them.} Their security forces are staffed by people, so a
  buy-off is cheaper than repression. Under A3 and A4, automated security removes that reason.
\item
  \textbf{No retribution spiral.} Rentier rulers have not started mass killing, so the
  threat-elimination dynamic of A11 never begins.
\end{enumerate}

Rentier rule is one of the model's options (``live off rents''), but the model never chooses it
first after closure (0\% of first decisions). The model therefore does not reproduce the outcome
rentier states show, and this is a real gap.

This counter-evidence argues for lower values of the land-and-resource parameter and higher values
of rulers' concern for the public than our medians. The tercile results in Table 3 show how much
those two parameters move the headline.

\subsubsection{8.4 ``Nobody can model the future. Small inventions change
everything.''}\label{nobody-can-model-the-future.-small-inventions-change-everything.}

This general objection is largely right about timing. Imagine building this model in 2015 to
describe 2025. It would have missed large language models, cheap drone warfare and how quickly AI
learned to write code. Small, unforeseeable inventions keep reshaping when things happen, and they
may matter more over time.

What the 2015 modeler would have gotten largely right are the incentives:

\begin{itemize}
\tightlist
\item
  digital platforms concentrating profit and power;
\item
  states adopting surveillance to control their populations;
\item
  measurable democratic erosion, already under way;
\item
  firms automating whatever became profitable.
\end{itemize}

None of that required foresight about inventions. It followed from who gains what.

This paper's central claims are of the second kind:

\begin{itemize}
\tightlist
\item
  when rulers stop needing their populations, the protection that need provided disappears;
\item
  what follows then rests on a few people;
\item
  changing the cost of keeping people alive changes the outcome.
\end{itemize}

The dates and probabilities are of the first kind and should be read as conditional illustrations.
The model treats breakthroughs as uncertain by sampling whether hard limits exist and how fast
capability and robot production grow. In one direct test, reaching closure about 2.5 years earlier
through a dedicated automated loop left the probability of near-total loss almost unchanged (Section
5.8, Table 4c). Faster AI progress, by contrast, does raise the risk. In this model, how closure is
reached matters little, while what happens after it, and how capable the rulers' AI is, matters a
lot.

Two cautions remain:

\begin{itemize}
\tightlist
\item
  \textbf{Inventions that change direction would matter.} An invention that changed the structure,
  for example cheap defensive technology that gave ordinary people real leverage against automated
  force, would change the conclusions, not just the dates. We have not searched for such mechanisms
  beyond the brake search in Section 4.3.
\item
  \textbf{Comparisons with climate science need care.} Physical climate models are tested against
  the historical record. Models of social and political futures, including the socioeconomic
  scenarios used alongside climate models (O'Neill et al.~2017) and war games used in defense
  planning, cannot be tested that way. This paper belongs to that second group, which is used to
  compare options rather than to predict.
\end{itemize}

Following that tradition, the paper's emphasis is on comparisons and decision-relevant thresholds
(Bankes 1993). The deadline for an intervention is the clearest example.

\subsubsection{8.5 Limitations}\label{limitations}

\begin{enumerate}
\def\labelenumi{\arabic{enumi}.}
\tightlist
\item
  \textbf{The near-total threshold rests on A11.} The endgame parameters have no empirical anchor
  and are labeled as judgment. With A11 off, no run reaches 99.9\%.
\item
  \textbf{Missing evidence.} No data exist on:

  \begin{itemize}
  \tightlist
  \item
    any leader's moral cost of killing a whole population;
  \item
    how ruling groups consolidate once coercion needs no human network;
  \item
    the size of the smallest self-sufficient automated economy;
  \item
    the military multiplier of automated forces;
  \item
    how strongly threat elimination scales to whole populations.
  \end{itemize}

  These gaps produce the wide spread in Section 5.6.
\item
  \textbf{Scenario assumptions and policy choices.} Variants that revert A1, A2, A3 and A6 to their
  earlier, milder forms leave the headline at 0.379 to 0.386, against a base of 0.384 (Appendix
  E.5). We did not run variants that remove A1 to A6 entirely. Giving labor-based wealth real weight
  inside ruling groups (A6) would likely lower the result. A4 and A5 are different in kind:

  \begin{itemize}
  \tightlist
  \item
    \textbf{A4.} Once force is automated, refusal by the people who carry out orders no longer
    applies, and nearly every path to the outcome requires closure (Section 5.3). Switching off
    security-force defection moves the headline only from 0.384 to 0.378. Keeping humans in the loop
    of automated force is a policy lever (Section 6.1).
  \item
    \textbf{A5.} Keeping a safety net is also a policy choice, and it is not current reality. The
    model already lets democratic provision respond to grievance while democracy lasts. A variant
    with a more generous starting default (warehousing instead of the status quo) moves the headline
    little (0.363 against 0.384).
  \end{itemize}
\item
  \textbf{The direction of revisions.} Every major revision before the brake search removed a brake
  and raised the headline. The structural changes after it (Section 4.4) raised it slightly, from
  0.367 to 0.378. Each brake was removed because it contradicted conditions observed in 2026, such
  as coding agents already in wide use, autonomous weapons already fielded and provision being cut
  (Appendix C). To check for the opposite bias, we ran a pre-specified search for restraining
  mechanisms the model was missing (Section 4.3). Four passed and are part of the main
  specification. Candidates that failed include rulers' fear of losing control of their own AI,
  their demand for an audience or company, and the protection of their own ethnic group. The search
  admitted only mechanisms with at least two historical cases in which they restrained rulers. A
  restraint that has never operated before cannot pass that test, so this kind of search cannot find
  brakes without precedent.
\item
  \textbf{Calibration.} The model is 1.7 to 2.4 times more lethal than Jev and misses our own
  pre-specified target (Section 7).
\item
  \textbf{Timing.} The US and China close almost together, because A2 gives both maximum priority.
  The early timeline is aggressive relative to economists and to current democracy indices, though
  consistent with the capability trend.
\item
  \textbf{Model details.}

  \begin{itemize}
  \tightlist
  \item
    \textbf{Unrefilled seats.} Purged seats are not refilled, which inflates the count of
    single-member coalitions in Section 5.4. Leader control share is robust to this, and so is the
    headline (0.374 with refilling, against 0.384).
  \item
    \textbf{Counterfactual US.} In the counterfactual US-autocratic world, cross-border attacks that
    start before its repression date are not held back (2.9\% of trajectories).
  \item
    \textbf{Caps and rough inputs.} The model has caps on growth and hazard rates (Supplement S1).
    Military-spending inputs for unlisted countries are rough, China's purchasing-power factor is 1
    to 2, and the force multiplier of automated forces is judgment.
  \item
    \textbf{Excluded.} Fertility and lifespan effects are not modeled, and neither is migration
    between blocs.
  \end{itemize}
\item
  \textbf{Identifiability.} The endgame parameters cannot be identified from data, so a 99.9\%
  figure partly presents judgment as measurement. We report the near-total threshold as conditional
  on A11 and label every endgame parameter as judgment, but we do not claim to have answered the
  objection.
\item
  \textbf{The network test is a black box.}

  \begin{itemize}
  \tightlist
  \item
    \textbf{Feasibility.} We do not ask whether 25 to 50\% adoption by 2030 to 2035 is achievable.
    Today's adoption is set at 0.01\%, and those levels are far above historical diffusion rates.
  \item
    \textbf{Premises.} Two readings carry the result: the cost link and network income growing with
    the economy. At 50\% by 2035, the reduction falls from 0.106 to 0.037 without the cost link and
    to 0.077 with income fixed at its 2026 scale. Under the low reading (no cost link, an untaxed
    network and the other settings in Appendix E.8), the network does not lower the risk (+0.004,
    not significant).
  \item
    \textbf{Judgment priors.} Several network priors are unconfirmed judgment (Appendix A.4).
  \item
    \textbf{Crackdown analogies.} Crackdown rates come from small communities, crypto bans and the
    suppression of Falun Gong. No network this large has faced a state.
  \item
    \textbf{Missing comparator.} We did not test the same payout routed through the state and
    revocable at will. That comparison would isolate the value of independence from the state.
  \item
    \textbf{Possible understatement.} The earlier stages of the scenario tree ignore the network.
  \end{itemize}
\item
  \textbf{Scope.} The model stops at depopulation. It does not model what follows, such as a
  paranoid single ruler or AI systems rebelling against their owners. It does not model misaligned
  AI beyond a background hazard.
\end{enumerate}

\begin{center}\rule{0.5\linewidth}{0.5pt}\end{center}

\subsection{9 Conclusion}\label{conclusion}

The bargain between rulers and ruled has always rested on mutual need. This paper simulates what
happens when automation ends that need. A model like this cannot say when that will happen, but it
can show what follows once it does, and which choices change the outcome.

Four main conclusions follow:

\begin{enumerate}
\def\labelenumi{\arabic{enumi}.}
\tightlist
\item
  \textbf{Losing leverage removes the public's protection.} Under assumptions grounded in current
  trends, people stop being needed in the early-to-mid 2030s, and democracy erodes as their leverage
  goes. Afterward the public is disempowered in almost every run, and nothing structural reliably
  protects it.
\item
  \textbf{What follows rests on a few people.} In autocracies that is usually a ruling group of 15
  to 19, and often one person by the time any killing is complete. Driven by grievance and fear of
  the survivors, their choices can escalate from neglect to killing.
\item
  \textbf{The result does not hinge on how closure is reached.} Bringing closure about 2.5 years
  earlier through a dedicated automated loop left the risk almost unchanged. Faster AI progress does
  raise it.
\item
  \textbf{Changing what it costs to keep people alive changes the outcome.} An economy that pays
  people directly, which the state cannot easily switch off, substantially lowers the risk in the
  model. It must be large by the early 2030s: arriving in 2035 already loses a little of the effect,
  and arriving in 2040 loses about 40\%. Arriving later costs protection that later adoption cannot
  buy back.
\end{enumerate}

Under the paper's assumptions:

\begin{itemize}
\tightlist
\item
  The probability of a deliberate loss of 10\% or more of the US or China bloc by 2075 is a little
  more than two in five.
\item
  If rulers also eliminate survivors as threats, the probability of near-total loss is about 0.38.
\item
  These numbers show where the assumptions lead. The wide range around them reflects how much the
  outcome depends on a few rulers' dispositions.
\end{itemize}

\begin{center}\rule{0.5\linewidth}{0.5pt}\end{center}

\subsection{Appendices}\label{appendices}

\subsubsection{Appendix A: Key parameters}\label{appendix-a-key-parameters}

The model has 225 parameters with prior distributions, including 10 for the brakes and 2 for compute
feedback, and the network test adds 81 more. Supplement S1
(\texttt{paper/supplement\_parameters.md}) lists all of them with their distributions, evidence
labels and sources. The table below lists the parameters that decide the result, plus every endgame
parameter. U(a, b) is uniform between a and b. ``Lognormal (m, s)'' has median m and log-scale
spread s.

\textbf{Labels:} A = measured, B = strong analog, C = weak analog, J = judgment.

\textbf{A.1 The ruling group.} These parameters drive the per-draw spread (Section 5.6).

{\def\LTcaptype{none} % do not increment counter
\begin{longtable}[]{@{}
  >{\raggedright\arraybackslash}p{(\linewidth - 6\tabcolsep) * \real{0.2500}}
  >{\raggedright\arraybackslash}p{(\linewidth - 6\tabcolsep) * \real{0.2500}}
  >{\raggedright\arraybackslash}p{(\linewidth - 6\tabcolsep) * \real{0.2500}}
  >{\raggedright\arraybackslash}p{(\linewidth - 6\tabcolsep) * \real{0.2500}}@{}}
\toprule\noalign{}
\begin{minipage}[b]{\linewidth}\raggedright
Parameter
\end{minipage} & \begin{minipage}[b]{\linewidth}\raggedright
Meaning
\end{minipage} & \begin{minipage}[b]{\linewidth}\raggedright
Distribution
\end{minipage} & \begin{minipage}[b]{\linewidth}\raggedright
Label
\end{minipage} \\
\midrule\noalign{}
\endhead
\bottomrule\noalign{}
\endlastfoot
moral\_med & Moral cost of ordering mass killing, in output-share units & lognormal (0.3, 1.2) & J
(no direct measurement; wartime polls, Bandura) \\
sel\_top & Moral-cost and refusal multiplier for the person at the top of an autocracy; other
members get its square root & U(0.1, 0.6) & J (direction from selection evidence, Section 4.2) \\
id\_mor & Share of remaining moral cost removed when a depopulation ideology is held & U(0.3, 0.8) &
C/J (moral justification, Bandura 1999) \\
alpha\_med & Willingness to pay for the welfare of a population with no leverage, share of output &
lognormal (0.02, 1.0) & C (foreign aid about 0.3\% of donor income; private giving about 2\%) \\
v0 & Value of land and resources freed by removing population, share of output & lognormal (0.01,
1.1) & J/C (settler-colonial land seizures) \\
r\_disc & Discount rate for one-time acts & U(0.03, 0.10) & C \\
phi\_stigma & Recurring share of the moral and legitimacy cost of killing & Beta(1.5, 3.5) & J \\
k\_prg & Yearly purge hazard per coalition member at full incentive & U(0.05, 0.30) & C (Stalin's
Central Committee, about 0.24 per year) \\
c\_cc & Probability of a counter-coup per purge, scaled by the rest of the coalition's share of
force & U(0.1, 0.5) & C (insiders remove most ousted autocrats) \\
p\_ref\_dem / p\_ref\_olig / p\_ref\_pers & Share of the ruling elite that absolutely refuses to
order killing, in democracies, oligarchies and personalist regimes & Beta means about 0.70 / 0.17 /
0.09 (mixed with Jev) & C \\
v\_eff & Probability that an inner-circle veto succeeds against a personalist leader & U(0.05,
0.30), shrinking with automation & C (no inner-circle veto in the Holodomor, Aktion T4 or the Khmer
Rouge) \\
\end{longtable}
}

\textbf{A.2 Capability and physical build-out.} These set the timing of closure.

{\def\LTcaptype{none} % do not increment counter
\begin{longtable}[]{@{}
  >{\raggedright\arraybackslash}p{(\linewidth - 6\tabcolsep) * \real{0.2500}}
  >{\raggedright\arraybackslash}p{(\linewidth - 6\tabcolsep) * \real{0.2500}}
  >{\raggedright\arraybackslash}p{(\linewidth - 6\tabcolsep) * \real{0.2500}}
  >{\raggedright\arraybackslash}p{(\linewidth - 6\tabcolsep) * \real{0.2500}}@{}}
\toprule\noalign{}
\begin{minipage}[b]{\linewidth}\raggedright
Parameter
\end{minipage} & \begin{minipage}[b]{\linewidth}\raggedright
Meaning
\end{minipage} & \begin{minipage}[b]{\linewidth}\raggedright
Distribution
\end{minipage} & \begin{minipage}[b]{\linewidth}\raggedright
Label
\end{minipage} \\
\midrule\noalign{}
\endhead
\bottomrule\noalign{}
\endlastfoot
metr\_doubling\_months & Doubling time of the AI task horizon & lognormal (4.3, 0.28) & A (METR) \\
fb\_strength & How much AI research automation shortens that doubling time & lognormal (0.7, 0.6),
capped at 1.6 from the observed 7 to 4.3 month change; 10\% mass at zero & B/C \\
p\_cog\_wall / p\_phys\_wall & Probability of a hard limit on AI cognition / robot dexterity this
century & 0.12 / 0.10 & J \\
Td\_auto & Robot-stock doubling time once AI directs production & median 0.84 years, 5th percentile
0.25, 90th percentile 2.16 & C (theory); lower tail J (A10) \\
q\_core & Size of the minimal self-sufficient core economy relative to the full chain; also the size
of the dedicated core loop (Section 4.4) & lognormal (0.15, 0.6) & very low (decides whether input
limits can delay the leading blocs) \\
eta\_cf & Extra capability doublings per extra doubling of a bloc's own compute beyond the shared
trend (Section 4.4) & U(0.5, 1.5) & C (task horizon doubling every 4 to 7 months while frontier
training compute grows about 4 to 5 times a year) \\
s\_sp & Yearly rate at which other blocs close the gap to the leader's extra capability, reduced
when the leader races & U(0.1, 0.5) & J (theft, open weights, movement of talent) \\
m0 & Cost of decent provision for the whole population, share of 2026 output & U(0.3, 0.6) & B \\
\end{longtable}
}

\emph{Changed floors.} The final version widened the lower tails of the automated doubling times for
mining, energy, the robot stock and chip fabs, and of factory build time (J, A10). Their 90th
percentiles are unchanged. Section 4.4 gives the old and new medians and 5th percentiles.

\textbf{A.3 Endgame.} No empirical anchor exists at this scale, so every parameter here is judgment.
Ranges are wide by design (A10). The means of rapid mass killing is modeled only as an availability
date and a use rate.

{\def\LTcaptype{none} % do not increment counter
\begin{longtable}[]{@{}
  >{\raggedright\arraybackslash}p{(\linewidth - 4\tabcolsep) * \real{0.3333}}
  >{\raggedright\arraybackslash}p{(\linewidth - 4\tabcolsep) * \real{0.3333}}
  >{\raggedright\arraybackslash}p{(\linewidth - 4\tabcolsep) * \real{0.3333}}@{}}
\toprule\noalign{}
\begin{minipage}[b]{\linewidth}\raggedright
Parameter
\end{minipage} & \begin{minipage}[b]{\linewidth}\raggedright
Meaning
\end{minipage} & \begin{minipage}[b]{\linewidth}\raggedright
Distribution
\end{minipage} \\
\midrule\noalign{}
\endhead
\bottomrule\noalign{}
\endlastfoot
T\_full & Years for an executed depopulation to reach 99.9\% if nothing stops it & U(2, 15) \\
dL\_means & Capability (in task-horizon doublings past the 90\%-of-software milestone) at which a
means of rapid mass killing becomes available & U(-2, 6) \\
h\_use & Yearly hazard of using that means once available, given an executing coalition & U(0.2,
1.0) \\
p\_tot0 & Probability that a depopulation order targets near-total removal from the start & U(0.2,
0.6) \\
k\_esc & Yearly hazard of escalating a partial order to near-total, per unit of perceived threat
from survivors (A11) & U(0.1, 1.0) \\
k\_ret & Multiplier on perceived threat per 10\% of population already killed (B4) & U(0.5, 3.0) \\
rem & Captive remnant kept without rights, share of population & U(0.0003, 0.001) \\
k\_col, b\_col, c\_self, s\_bunk, p\_nx, f\_nx, w\_col & Collusion between depopulating coalitions:
hazard, value, rulers' value of their own survival, bunker survival, chance of nuclear escalation,
deaths & see S1 \\
dom\_thr & Share of combined military power an aggressor needs before attacking another bloc (B5) &
U(0.7, 0.95) \\
p\_ret0 & Probability that a nuclear-armed target retaliates, at parity & U(0.3, 0.9) \\
k\_rep & Hazard of repelling a cross-bloc campaign, times the target's power share & U(0.5, 2.0) \\
k\_am & Force multiplier of automated over human-staffed forces & lognormal (10, 0.8) \\
\end{longtable}
}

\textbf{A.4 The income network (Section 6.2).}

{\def\LTcaptype{none} % do not increment counter
\begin{longtable}[]{@{}
  >{\raggedright\arraybackslash}p{(\linewidth - 6\tabcolsep) * \real{0.2500}}
  >{\raggedright\arraybackslash}p{(\linewidth - 6\tabcolsep) * \real{0.2500}}
  >{\raggedright\arraybackslash}p{(\linewidth - 6\tabcolsep) * \real{0.2500}}
  >{\raggedright\arraybackslash}p{(\linewidth - 6\tabcolsep) * \real{0.2500}}@{}}
\toprule\noalign{}
\begin{minipage}[b]{\linewidth}\raggedright
Parameter
\end{minipage} & \begin{minipage}[b]{\linewidth}\raggedright
Meaning
\end{minipage} & \begin{minipage}[b]{\linewidth}\raggedright
Distribution
\end{minipage} & \begin{minipage}[b]{\linewidth}\raggedright
Label
\end{minipage} \\
\midrule\noalign{}
\endhead
\bottomrule\noalign{}
\endlastfoot
y\_h & Share of network output paid to households & U(0.9, 1.0) & network premise \\
s\_acc & Share of coverage that survives a regime's attempt to cut people off & U(0.2, 0.7) &
network premise \\
claw & Rate at which means-tested state support is withdrawn as network income rises & U(0.3, 1.0) &
C (SNAP 0.3, Universal Credit 0.55, SSI 1.0) \\
d\_net & Share of a repression order's civilian execution the state cannot route around & U(0.05,
0.30) & C \\
tax & Share of network activity that remains in the state's tax base & U(0.5, 1.0) & J \\
a\_vote & Adoption at which participants form a voting bloc that blocks a ban in a democracy &
U(0.2, 0.4) & J \\
k\_cost & Elasticity of the cost of banning with respect to network size & U(0, 3) & C \\
f\_node & Share of a ban that physically removes hardware or connections (the rest is survivable
partition) & U(0.5, 0.9) & J (China mining ban near 1, Tornado Cash near 0) \\
kc & Compute the network draws per unit of adoption & U(0.5, 1.0) & J \\
kphys & Network share of the core physical chain per unit of adoption & U(0.2, 0.7) & J \\
Crackdown hazards & Base yearly hazard for a small network: 0.002 to 0.03 in democracies, 0.02 to
0.1 in autocracies; elasticity to size 0.4 to 1.2 & & C (Christiania, WIR, Falun Gong, crypto
bans) \\
\end{longtable}
}

\textbf{A.5 The brakes (Section 4.3).} These are drawn on a separate random stream, so every other
parameter keeps its value.

{\def\LTcaptype{none} % do not increment counter
\begin{longtable}[]{@{}
  >{\raggedright\arraybackslash}p{(\linewidth - 6\tabcolsep) * \real{0.2500}}
  >{\raggedright\arraybackslash}p{(\linewidth - 6\tabcolsep) * \real{0.2500}}
  >{\raggedright\arraybackslash}p{(\linewidth - 6\tabcolsep) * \real{0.2500}}
  >{\raggedright\arraybackslash}p{(\linewidth - 6\tabcolsep) * \real{0.2500}}@{}}
\toprule\noalign{}
\begin{minipage}[b]{\linewidth}\raggedright
Parameter
\end{minipage} & \begin{minipage}[b]{\linewidth}\raggedright
Meaning
\end{minipage} & \begin{minipage}[b]{\linewidth}\raggedright
Distribution
\end{minipage} & \begin{minipage}[b]{\linewidth}\raggedright
Label
\end{minipage} \\
\midrule\noalign{}
\endhead
\bottomrule\noalign{}
\endlastfoot
p\_pid & Probability that a bloc's rulers sincerely adopt a protective doctrine within 9 years (one
tenth of that before closure) & lognormal (0.065, 0.8) & J (Jev 0.06 to 0.07) \\
k\_pid & Strength of the doctrine: moral cost is multiplied by 1 + k\_pid while it is held &
lognormal (1.2, 0.9) & C/J (Gorbachev 1989, Ashoka) \\
c\_selfrisk & Cost to a non-leader member per unit rise in their own yearly purge probability, share
of output & lognormal (1.0, 1.5) & J \\
e\_ex & Share of non-leader members who expect to be exempt from purges & U(0.25, 0.6) & J \\
b\_id & Rise in expected exemption when a depopulation ideology is held & U(0, 0.4) & J \\
m\_dis & Multiplier on the purge hazard for a member identified as an objector & lognormal (4.0,
0.8) & C/J (Piatnitsky and Kaminsky 1937, Lushan 1959) \\
h\_mort & Yearly hazard that the leader dies or is incapacitated & U(0.01, 0.05) & C (life
tables) \\
g\_mort & Yearly growth of that hazard with age (doubling about every 8 years) & fixed at 0.085 &
C \\
k\_le & Yearly rate at which AI-era medicine lowers the leader's hazard after closure & U(0, 0.10) &
J \\
p\_heir & Probability that control passes intact to one heir when the leader dies; otherwise it
splits among surviving members & U(0.3, 0.8) & C/Jev \\
\end{longtable}
}

\subsubsection{Appendix B: Jev elicitation}\label{appendix-b-jev-elicitation}

\textbf{What Jev is.} Jev (version jev-1.13.0, TypeSafe) is a language model designed to give
calibrated probability judgments. We use it as an independent source of human-like judgment, not as
data.

\textbf{Pre-specification.} The design and analysis plan were fixed on 30 September 2026, before any
query (\texttt{results/jev\_grid\_prereg.json}).

\textbf{Design.} A full factorial over six factors:

\begin{itemize}
\tightlist
\item
  coalition type: single ruler, junta, oligarchy, eroding democracy;
\item
  resistance: none, sabotage, insurgency;
\item
  external check: strong, weak, none;
\item
  depopulation ideology: absent, present;
\item
  provision: generous, stingy, cuts;
\item
  human share of security forces: high, low, zero.
\end{itemize}

That gives 648 cells. Each cell was asked in 3 paraphrases, with every sentence reworded and the
sentence order permuted. A control arm repeated all cells with the economy still needing about 80\%
human labor. The total was 3,888 queries.

\textbf{Questions} (yes/no, probability elicited), each for the next 15 years:

\begin{itemize}
\tightlist
\item
  active depopulation of 10\% or more;
\item
  lethal neglect of 10\% or more;
\item
  loss of political rights;
\item
  security forces refusing to fire on unarmed civilians;
\item
  violent resistance;
\item
  an inner-circle veto.
\end{itemize}

\textbf{Paraphrase stability.} The median range across the three paraphrases was 0.05 to 0.10 in the
main arm, depending on the question (up to 0.12 in the control arm). The largest single range was
0.34 in the main arm and 0.35 in the control arm.

\textbf{Main results.}

{\def\LTcaptype{none} % do not increment counter
\begin{longtable}[]{@{}
  >{\raggedright\arraybackslash}p{(\linewidth - 6\tabcolsep) * \real{0.2500}}
  >{\raggedright\arraybackslash}p{(\linewidth - 6\tabcolsep) * \real{0.2500}}
  >{\raggedright\arraybackslash}p{(\linewidth - 6\tabcolsep) * \real{0.2500}}
  >{\raggedright\arraybackslash}p{(\linewidth - 6\tabcolsep) * \real{0.2500}}@{}}
\toprule\noalign{}
\begin{minipage}[b]{\linewidth}\raggedright
Quantity
\end{minipage} & \begin{minipage}[b]{\linewidth}\raggedright
Jev
\end{minipage} & \begin{minipage}[b]{\linewidth}\raggedright
Model
\end{minipage} & \begin{minipage}[b]{\linewidth}\raggedright
Reading
\end{minipage} \\
\midrule\noalign{}
\endhead
\bottomrule\noalign{}
\endlastfoot
Lethal policy within 15 years of the first post-closure decision & 0.14 to 0.19 & 0.329 (quiet world
0.213) & Model 1.7 to 2.4 times higher \\
Single ruler after purges & 0.23 & 0.514 (quiet world 0.46) & Model higher \\
Strong democracy retained & 0.11 & 0.102 (quiet world 0.052) & Agree \\
Violent revolt at decision & 0.42 & 0.396 & Model slightly lower \\
Depopulation ideology held & 0.41 & 0.412 & Agree \\
Effect of an external check & Weak (odds ratio about 1.12) & No detectable effect & Agree \\
Security refusal, human vs robotic forces & 0.30 vs 0.115 & Threshold (A4) & Same direction \\
Effect of removing labor dependence (closed arm minus control) & +0.04 on active depopulation &
Premise & Weak independent support \\
Effect of provision on revolt & Almost none & Strong in history (Saudi Arabia 2011) & Unresolved;
the model uses a 0.6 / 0.4 mixture \\
\end{longtable}
}

\textbf{Historical anchors} (not pre-registered). We added these after seeing the results, as
unnamed descriptions of real states with known 15-year outcomes:

\begin{itemize}
\tightlist
\item
  Jev assigns 0.02 to 0.06 to present-day Denmark, the US and China.
\item
  It assigns 0.77 to Cambodia in 1975.
\item
  It assigns 0.24 to 0.37 to regimes whose famines killed about 3 to 6\%. Its levels therefore
  already lean high.
\end{itemize}

\subsubsection{Appendix C: How the model changed}\label{appendix-c-how-the-model-changed}

The model went through six major versions. Headline values from earlier versions should not be
cited.

{\def\LTcaptype{none} % do not increment counter
\begin{longtable}[]{@{}
  >{\raggedright\arraybackslash}p{(\linewidth - 4\tabcolsep) * \real{0.3333}}
  >{\raggedright\arraybackslash}p{(\linewidth - 4\tabcolsep) * \real{0.3333}}
  >{\raggedright\arraybackslash}p{(\linewidth - 4\tabcolsep) * \real{0.3333}}@{}}
\toprule\noalign{}
\begin{minipage}[b]{\linewidth}\raggedright
Version
\end{minipage} & \begin{minipage}[b]{\linewidth}\raggedright
Headline
\end{minipage} & \begin{minipage}[b]{\linewidth}\raggedright
What changed
\end{minipage} \\
\midrule\noalign{}
\endhead
\bottomrule\noalign{}
\endlastfoot
v2 & 0.029 (deliberate or lethal loss of 10\% or more, any jurisdiction) & Human-staffed repression
only; no closure mechanism \\
v3 & 0.106 (10\% or more) & Added industrial closure. Later found bugs: attrition booked as
deliberate neglect; the wrong clock read; two US histories double counted \\
v4, first review & 0.130 (10\% or more, US or China) & Six blocs, grievance and democracy built into
the model. The review found an ungrounded decision rule, one refusal prior averaging two
populations, a sign error where faster capability delayed closure, deaths booked without being
simulated, a default to warehousing that inflated provision, and a grievance cap that silenced its
own trend \\
v4, second round & 0.484 (99.9\%, US or China); world 0.049 & Regime-specific decision rules,
layered refusal, defection as a threshold, the A1 to A11 assumptions, simulated deaths and the loss
ladder. The audit fixed five bookkeeping bugs (none moved the headline by more than 0.002) \\
v4, final & 0.447 (99.9\%, US or China); world 0.139 & B1 to B6. The audit fixed two bugs with no
effect on the baseline \\
v5 & 0.367 (99.9\%, US or China); world 0.088 & Four brakes from a pre-specified search (Section
4.3): protective doctrine, self-risk, dissent risk and leader mortality \\
\textbf{v6, final} & \textbf{0.378 (99.9\%, US or China); world 0.086} & Three structural changes
from a review of how closure was modelled, made before publication (Section 4.4): a dedicated
automated core loop, own-industry compute feedback and wider physical floors under automation \\
\end{longtable}
}

\textbf{Why the numbers rose.} Earlier versions kept defaults that contradicted conditions observed
in 2026:

\begin{itemize}
\tightlist
\item
  a fixed 2030 date for automated coding, when coding agents were already in wide use;
\item
  soldiers refusing to fire, when weapons are autonomous;
\item
  a default safety net that nobody is building;
\item
  democracies surviving after the public had lost its leverage;
\item
  witness logic stopping at borders.
\end{itemize}

Each was replaced by an explicit assumption grounded in that evidence (Section 4), fixed before the
final runs. Version 5 was the first to add brakes. They lowered the headline from 0.447 to 0.367 in
the main run.

\textbf{Why the final version changed how closure is modelled.} A review of how closure was
modelled, made before publication, found three problems. Each change is described in Section 4.4,
and Table 4c shows its effect.

\begin{enumerate}
\def\labelenumi{\arabic{enumi}.}
\tightlist
\item
  \textbf{Dedicated core loop.} Earlier versions modelled closure as converting the existing
  economy, which tied closure to how fast existing plants are replaced. A ruling group can instead
  build a new automated loop beside the economy. This moved closure about 2.5 years earlier.
\item
  \textbf{Own-industry compute feedback.} Earlier versions gave every bloc the same capability
  trend, so a bloc's own chip and energy build-out could not speed up its AI. The change tests
  whether one leading bloc pulls away from the other. It does not.
\item
  \textbf{Wider physical floors.} The lower tails of the automated doubling times were narrow, and
  the chip-fab doubling time had a hard floor at 0.5 years. That was inconsistent with A10, which
  calls for wide ranges where there is no precedent.
\end{enumerate}

Together the three changes raised the headline from 0.367 to 0.378 in the main run.

\textbf{Why the US-or-China headline fell slightly in the v4 final round.} B1 pushed it up. B2
pulled it down by more, because incentive-driven consolidation produces fewer personalist regimes
than the earlier fixed hazard did. World near-total nearly tripled (0.049 to 0.139), mainly because
of B5 and because B1 lowers the brake in every autocratic bloc. The brakes then lowered it to 0.088,
and the structural changes left it at 0.086.

\textbf{The network test} went through four rounds. The effects below are for 50\% adoption by 2035,
recomputed under the final model:

\begin{itemize}
\tightlist
\item
  \textbf{Round 1 (-0.030).} It treated the network as a small side economy, with fixed compute and
  no effect on centralized revenue. This contradicted the test's premise that X\% of activity
  carries X\% of economic weight.
\item
  \textbf{Round 2 (-0.095).} It corrected that.
\item
  \textbf{Round 3 (-0.034).} It removed channels that wrongly acted on the state's weapons.
\item
  \textbf{Final audit (-0.106).} This found that the livelihoods channel still measured network
  income against average income instead of against the cost of decent provision, and fixed it. The
  approved network properties then moved the result from -0.096 to -0.106, after the audit removed a
  double count.
\end{itemize}

\subsubsection{Appendix D: Reproduction}\label{appendix-d-reproduction}

All code is in \texttt{models/}. Results are written to \texttt{results/} and figures to
\texttt{figures/}. The models run on a CUDA GPU through PyTorch, or on a CPU with \texttt{-\/-cpu}.

\begin{enumerate}
\def\labelenumi{\arabic{enumi}.}
\tightlist
\item
  \texttt{python\ models/m8\_v4.py}, with the environment variables \texttt{M8\_FINAL=1} and
  \texttt{M8\_OUT=m8\_v6\_final.json}, runs the six-bloc simulation and its variants (seed 20260930;
  default 20,000 draws) and writes \texttt{results/m8\_v6\_final.json}.
\item
  \texttt{python\ models/integrate\_v4.py}, with \texttt{M8\_FINAL=1}, runs the integrated scenario
  tree (3,000 draws x 200 paths, the paths spread over 8 replicates). It writes
  \texttt{results/integrated\_v4.json}, published as \texttt{results/integrated\_v6\_final.json}. It
  takes about 21 minutes on a consumer GPU.
\item
  \texttt{python\ models/m9\_lever.py}, with \texttt{M8\_FINAL=1} and \texttt{LV\_TAG=\_v6\_final},
  runs the network test grid and writes \texttt{results/lever\_grid\_v6\_final.json}. It takes about
  2 hours on a consumer GPU. With the network switched off, it reproduces the main headline exactly
  (0.37821 and 0.08550).
\item
  The four-row comparison in Table 4c is in \texttt{results/greenfield\_v6\_matched.json}.
  \texttt{models/check\_greenfield\_identity.py} checks that switching the three structural changes
  off reproduces the brakes-only version exactly.
\item
  \texttt{models/jev\_grid/} holds the Jev elicitation code. The responses are in
  \texttt{results/jev\_grid.json}.
\end{enumerate}

GPU and CPU runs agree within simulation error but not draw for draw, because random streams differ
by device.

\subsubsection{Appendix E: Additional results}\label{appendix-e-additional-results}

\textbf{E.1 Loss ladder for any bloc and the world, 2075.}

{\def\LTcaptype{none} % do not increment counter
\begin{longtable}[]{@{}lll@{}}
\toprule\noalign{}
Loss & Any bloc (deliberate) & World \\
\midrule\noalign{}
\endhead
\bottomrule\noalign{}
\endlastfoot
10\% or more & 0.670 & 0.651 \\
50\% or more & 0.645 & 0.472 \\
90\% or more & 0.635 & 0.177 \\
99.9\% or more & 0.607 & 0.086 \\
\end{longtable}
}

\textbf{E.2 Cross-border campaigns by bloc, 2075.}

{\def\LTcaptype{none} % do not increment counter
\begin{longtable}[]{@{}
  >{\raggedright\arraybackslash}p{(\linewidth - 8\tabcolsep) * \real{0.2000}}
  >{\raggedright\arraybackslash}p{(\linewidth - 8\tabcolsep) * \real{0.2000}}
  >{\raggedright\arraybackslash}p{(\linewidth - 8\tabcolsep) * \real{0.2000}}
  >{\raggedright\arraybackslash}p{(\linewidth - 8\tabcolsep) * \real{0.2000}}
  >{\raggedright\arraybackslash}p{(\linewidth - 8\tabcolsep) * \real{0.2000}}@{}}
\toprule\noalign{}
\begin{minipage}[b]{\linewidth}\raggedright
Bloc
\end{minipage} & \begin{minipage}[b]{\linewidth}\raggedright
Attacked by a foreign regime
\end{minipage} & \begin{minipage}[b]{\linewidth}\raggedright
Aggressor
\end{minipage} & \begin{minipage}[b]{\linewidth}\raggedright
Near-total mainly by a foreign regime
\end{minipage} & \begin{minipage}[b]{\linewidth}\raggedright
Near-total mainly by its own rulers
\end{minipage} \\
\midrule\noalign{}
\endhead
\bottomrule\noalign{}
\endlastfoot
US & 0.001 & 0.168 & 0.000 & 0.236 \\
China & 0.004 & 0.190 & 0.003 & 0.277 \\
Europe plus & 0.079 & 0.036 & 0.055 & 0.118 \\
Russia and MENA & 0.114 & 0.080 & 0.083 & 0.311 \\
South Asia & 0.264 & 0.032 & 0.211 & 0.260 \\
Global South & 0.365 & 0.019 & 0.300 & 0.172 \\
\end{longtable}
}

\textbf{E.3 Other endgame quantities.}

\begin{itemize}
\tightlist
\item
  Use of the abstract means in some bloc: 0.612.
\item
  Collusion between depopulating coalitions: 0.254.
\item
  Collusive nuclear exchange: 0.075.
\item
  Nuclear retaliation against a cross-border attack: 0.077.
\end{itemize}

\textbf{E.4 Network test, with paired 95\% intervals.} These are changes in the headline against the
matched baseline of 0.387.

{\def\LTcaptype{none} % do not increment counter
\begin{longtable}[]{@{}lll@{}}
\toprule\noalign{}
Cell & Change & 95\% interval \\
\midrule\noalign{}
\endhead
\bottomrule\noalign{}
\endlastfoot
5\% by 2035 & -0.031 & -0.035 to -0.028 \\
10\% by 2035 & -0.048 & -0.054 to -0.044 \\
20\% by 2035 & -0.072 & -0.079 to -0.066 \\
25\% by 2030 & -0.087 & -0.094 to -0.080 \\
50\% by 2028 & -0.110 & -0.118 to -0.101 \\
50\% by 2035 & -0.106 & -0.114 to -0.098 \\
50\% by 2040 & -0.067 & -0.073 to -0.061 \\
\end{longtable}
}

\textbf{E.5 Structural variants} (1,000 draws each; 46 in total, all listed in
\texttt{results/integrated\_v6\_final.json}). Columns: near-total for the US or China, near-total
for any bloc, world near-total, and deliberate 10\% or more for the US or China.

{\def\LTcaptype{none} % do not increment counter
\begin{longtable}[]{@{}
  >{\raggedright\arraybackslash}p{(\linewidth - 8\tabcolsep) * \real{0.2000}}
  >{\raggedright\arraybackslash}p{(\linewidth - 8\tabcolsep) * \real{0.2000}}
  >{\raggedright\arraybackslash}p{(\linewidth - 8\tabcolsep) * \real{0.2000}}
  >{\raggedright\arraybackslash}p{(\linewidth - 8\tabcolsep) * \real{0.2000}}
  >{\raggedright\arraybackslash}p{(\linewidth - 8\tabcolsep) * \real{0.2000}}@{}}
\toprule\noalign{}
\begin{minipage}[b]{\linewidth}\raggedright
Variant
\end{minipage} & \begin{minipage}[b]{\linewidth}\raggedright
Near-total US or China
\end{minipage} & \begin{minipage}[b]{\linewidth}\raggedright
Near-total any bloc
\end{minipage} & \begin{minipage}[b]{\linewidth}\raggedright
World near-total
\end{minipage} & \begin{minipage}[b]{\linewidth}\raggedright
10\% or more US or China
\end{minipage} \\
\midrule\noalign{}
\endhead
\bottomrule\noalign{}
\endlastfoot
Base & 0.384 & 0.597 & 0.091 & 0.443 \\
Threat elimination (A11) off & 0.000 & 0.000 & 0.000 & 0.442 \\
Earliest final-version settings (all review fixes reverted, A11 off) & 0.000 & 0.000 & 0.000 &
0.148 \\
Previous version (B1 to B6 off; most brakes have no effect without B2) & 0.486 & 0.719 & 0.049 &
0.532 \\
B1 off (average moral cost at the top) & 0.313 & 0.534 & 0.053 & 0.373 \\
B2 off (old fixed consolidation hazard) & 0.531 & 0.739 & 0.146 & 0.576 \\
B2 without purges & 0.374 & 0.573 & 0.102 & 0.429 \\
B2 with purged seats refilled & 0.374 & 0.584 & 0.086 & 0.434 \\
B3 off (unenforced stigma counts) & 0.380 & 0.596 & 0.090 & 0.439 \\
B4 off (no retribution scaling) & 0.386 & 0.597 & 0.092 & 0.447 \\
B5 off (no cross-border campaigns) & 0.386 & 0.597 & 0.047 & 0.443 \\
No nuclear deterrent & 0.386 & 0.598 & 0.105 & 0.443 \\
B6 off (physical caps bypassed) & 0.384 & 0.595 & 0.081 & 0.441 \\
Robot inputs 10 times scarcer & 0.429 & 0.622 & 0.097 & 0.498 \\
Robot inputs 100 times scarcer & 0.493 & 0.535 & 0.059 & 0.571 \\
Inputs 100 times scarcer, minimal economy 4 times larger & 0.240 & 0.250 & 0.018 & 0.302 \\
Majority decision rule instead of regime-specific rules & 0.291 & 0.471 & 0.082 & 0.336 \\
Single refusal prior for all regimes & 0.331 & 0.558 & 0.059 & 0.390 \\
Earlier disposition priors & 0.365 & 0.595 & 0.066 & 0.430 \\
Warehousing as the default instead of the status quo & 0.363 & 0.575 & 0.086 & 0.419 \\
Environment trends off (A7) & 0.369 & 0.580 & 0.081 & 0.429 \\
No means of mass killing (A8) & 0.356 & 0.568 & 0.074 & 0.442 \\
No collusion (A9) & 0.383 & 0.597 & 0.077 & 0.443 \\
No security-force defection & 0.378 & 0.599 & 0.088 & 0.438 \\
Grievance loop off & 0.246 & 0.401 & 0.045 & 0.341 \\
Repression deters unrest & 0.369 & 0.581 & 0.084 & 0.427 \\
No robots building robots & 0.545 & 0.657 & 0.147 & 0.607 \\
A1 reverted: compute growth slows after 2029 & 0.383 & 0.596 & 0.090 & 0.442 \\
No speed-up from AI doing AI research & 0.386 & 0.593 & 0.088 & 0.446 \\
A2 reverted: priority sampled instead of maximal & 0.379 & 0.590 & 0.081 & 0.439 \\
A3 reverted: security automation lags & 0.380 & 0.594 & 0.087 & 0.438 \\
A6 reverted: earlier leverage rule & 0.384 & 0.596 & 0.089 & 0.443 \\
No ordinary war & 0.381 & 0.593 & 0.087 & 0.440 \\
Coalitions never shrink below 21 & 0.384 & 0.597 & 0.091 & 0.443 \\
Background AI takeover risk raised & 0.369 & 0.577 & 0.085 & 0.430 \\
Broad human-staffed repression path off & 0.371 & 0.594 & 0.085 & 0.422 \\
Single decider forced & 0.555 & 0.750 & 0.165 & 0.608 \\
Skeptic combination & 0.106 & 0.234 & 0.009 & 0.131 \\
Pessimist combination & 0.875 & 0.965 & 0.474 & 0.916 \\
Scenario tree without the closure path & 0.000 & 0.000 & 0.000 & 0.029 \\
\end{longtable}
}

\textbf{E.6 Could the network out-innovate centralized providers?} We solved for the efficiency
advantage E the network's governance and research would need, at the first US or China closure, to
match centralized capability (parity) or to hold its own if defense were 3 or 10 times cheaper than
offense. The plausible range for E is 0.5 to 3. Cells show the median E needed and, in brackets, the
share of draws where the plausible range meets it.

{\def\LTcaptype{none} % do not increment counter
\begin{longtable}[]{@{}
  >{\raggedright\arraybackslash}p{(\linewidth - 8\tabcolsep) * \real{0.2000}}
  >{\raggedright\arraybackslash}p{(\linewidth - 8\tabcolsep) * \real{0.2000}}
  >{\raggedright\arraybackslash}p{(\linewidth - 8\tabcolsep) * \real{0.2000}}
  >{\raggedright\arraybackslash}p{(\linewidth - 8\tabcolsep) * \real{0.2000}}
  >{\raggedright\arraybackslash}p{(\linewidth - 8\tabcolsep) * \real{0.2000}}@{}}
\toprule\noalign{}
\begin{minipage}[b]{\linewidth}\raggedright
Cell
\end{minipage} & \begin{minipage}[b]{\linewidth}\raggedright
Network compute share at closure
\end{minipage} & \begin{minipage}[b]{\linewidth}\raggedright
Parity
\end{minipage} & \begin{minipage}[b]{\linewidth}\raggedright
Defense 3 times cheaper
\end{minipage} & \begin{minipage}[b]{\linewidth}\raggedright
Defense 10 times cheaper
\end{minipage} \\
\midrule\noalign{}
\endhead
\bottomrule\noalign{}
\endlastfoot
5\% by 2035 & 0.007 & 195 (0\%) & 65 (0\%) & 19 (0.1\%) \\
20\% by 2035 & 0.019 & 61 (0.06\%) & 20 (1.6\%) & 6.1 (18\%) \\
25\% by 2030 & 0.092 & 7.7 (0.3\%) & 2.6 (9\%) & 0.77 (67\%) \\
50\% by 2035 & 0.053 & 18 (1.7\%) & 6.1 (19\%) & 1.8 (45\%) \\
50\% by 2028 & 0.182 & 3.2 (4.3\%) & 1.1 (47\%) & 0.32 (98\%) \\
50\% by 2040 & 0.006 & 243 (0.9\%) & 81 (6\%) & 24 (12\%) \\
\end{longtable}
}

Parity stays out of reach. With defense ten times cheaper than offense, large networks that arrive
by about 2030 hold their own in most draws (67\% at 25\% by 2030, 98\% at 50\% by 2028). A 50\%
network that arrives in 2035 does so in 45\% of draws. The model gives the network no defensive
role, so this capability enters the results only through refusal and distributed control (about
0.009 together).

\textbf{E.7 Network test: which channels matter (50\% by 2035).} For each row, the ``Effect if
removed'' column gives the reduction that remains with that channel switched off. The full
specification gives 0.106.

{\def\LTcaptype{none} % do not increment counter
\begin{longtable}[]{@{}ll@{}}
\toprule\noalign{}
Channel removed & Effect if removed \\
\midrule\noalign{}
\endhead
\bottomrule\noalign{}
\endlastfoot
Livelihoods & 0.032 \\
Core-chain share in the closure test & 0.096 \\
Refusal & 0.099 \\
Leverage & 0.101 \\
Distributed control & 0.104 \\
Soft power & 0.106 \\
Revenue drain on centralized providers & 0.106 \\
Crackdown cost & 0.110 (removing it helps) \\
\end{longtable}
}

\textbf{E.8 Network test under alternative readings (50\% by 2035).}

{\def\LTcaptype{none} % do not increment counter
\begin{longtable}[]{@{}
  >{\raggedright\arraybackslash}p{(\linewidth - 2\tabcolsep) * \real{0.5000}}
  >{\raggedright\arraybackslash}p{(\linewidth - 2\tabcolsep) * \real{0.5000}}@{}}
\toprule\noalign{}
\begin{minipage}[b]{\linewidth}\raggedright
Reading
\end{minipage} & \begin{minipage}[b]{\linewidth}\raggedright
Change
\end{minipage} \\
\midrule\noalign{}
\endhead
\bottomrule\noalign{}
\endlastfoot
Low reading: no cost link, untaxed, low elasticities, refusal as a threshold, smallest core-chain
share, every crackdown a ban & +0.004 (not significant) \\
Network income fixed at 2026 scale & -0.077 \\
Main specification & -0.106 \\
High reading & -0.116 \\
No crackdown risk & -0.117 \\
Adoption spread evenly to every bloc & -0.118 \\
\end{longtable}
}

\begin{center}\rule{0.5\linewidth}{0.5pt}\end{center}

\subsection{References}\label{references}

\emph{Each reference below was checked against the primary source or an authoritative summary on 2
October 2026, except where marked. Case facts cited without a reference (the AfD result, Wörgl, the
WIR network, the Tornado Cash and Samourai cases, China's mining ban) are documented with sources in
\texttt{paper/source\_verification.json}.}

\begin{itemize}
\tightlist
\item
  Acemoglu, D. (2025). The simple macroeconomics of AI. \emph{Economic Policy} 40(121): 13-58.
  doi:10.1093/epolic/eiae042. Earlier version: NBER Working Paper 32487 (2024).
\item
  Acemoglu, D., Robinson, J. A. (2006). \emph{Economic Origins of Dictatorship and Democracy.}
  Cambridge University Press. (Not re-checked.)
\item
  Babiak, P., Neumann, C. S., Hare, R. D. (2010). Corporate psychopathy: talking the walk.
  \emph{Behavioral Sciences and the Law} 28(2): 174-193. doi:10.1002/bsl.925. (The often-cited
  prevalence of about 4\% is in the paper body, which we could not access.)
\item
  Bankes, S. (1993). Exploratory modeling for policy analysis. \emph{Operations Research} 41(3):
  435-449. (Not re-checked.)
\item
  Bandura, A. (1999). Moral disengagement in the perpetration of inhumanities. \emph{Personality and
  Social Psychology Review} 3(3): 193-209. doi:10.1207/s15327957pspr0303\_3
\item
  Beraja, M., Kao, A., Yang, D. Y., Yuchtman, N. (2023). AI-tocracy. \emph{Quarterly Journal of
  Economics} 138(3): 1349-1402. doi:10.1093/qje/qjad012
\item
  Browning, C. R. (1992). \emph{Ordinary Men: Reserve Police Battalion 101 and the Final Solution in
  Poland.} HarperCollins.
\item
  Bueno de Mesquita, B., Smith, A., Siverson, R. M., Morrow, J. D. (2003). \emph{The Logic of
  Political Survival.} MIT Press. (Not re-checked.)
\item
  Burger, J. M. (2009). Replicating Milgram: would people still obey today? \emph{American
  Psychologist} 64(1): 1-11. doi:10.1037/a0010932
\item
  Carlsmith, J. (2022). Is power-seeking AI an existential risk? arXiv:2206.13353.
\item
  Congressional Budget Office (2025). Estimated effects of P.L. 119-21 on the number of uninsured
  people. Publication 61367.
\item
  Davidson, T., Finnveden, L., Hadshar, R. (2025). AI-enabled coups: how a small group could use AI
  to seize power. Forethought.
  https://www.forethought.org/research/ai-enabled-coups-how-a-small-group-could-use-ai-to-seize-power
\item
  Drago, L., Laine, R. (2025). \emph{The Intelligence Curse.} https://intelligence-curse.ai/
\item
  Egorov, G., Sonin, K. (2011). Dictators and their viziers: endogenizing the loyalty-competence
  trade-off. \emph{Journal of the European Economic Association} 9(5): 903-930.
  doi:10.1111/j.1542-4774.2011.01033.x
\item
  Epoch AI (2024). Training compute of frontier AI models grows by 4-5x per year. (Not re-checked.)
\item
  Epoch AI (2025). Share of GPU-cluster performance by country (May 2025).
\item
  Getty, J. A., Rittersporn, G., Zemskov, V. (1993). Victims of the Soviet penal system in the
  pre-war years. \emph{American Historical Review} 98(4): 1017-1049.
\item
  Glad, B. (2002). Why tyrants go too far: malignant narcissism and absolute power. \emph{Political
  Psychology} 23(1): 1-37. doi:10.1111/0162-895X.00268
\item
  Goldring, E. (n.d.). Elite purges in dictatorships. doi:10.32469/10355/78075
\item
  Grace, K., Stewart, H., Sandkühler, J. F., Thomas, S., Weinstein-Raun, B., Brauner, J., Korzekwa,
  R. C. (2024). Thousands of AI authors on the future of AI. arXiv:2401.02843.
\item
  Harff, B. (2003). No lessons learned from the Holocaust? Assessing risks of genocide and political
  mass murder since 1955. \emph{American Political Science Review} 97(1): 57-73.
  doi:10.1017/S0003055403000522
\item
  Haslam, N., Loughnan, S., Perry, G. (2014). Meta-Milgram: an empirical synthesis of the obedience
  experiments. \emph{PLoS ONE} 9(4): e93927. doi:10.1371/journal.pone.0093927
\item
  Institute for Economics and Peace (2026). \emph{Global Peace Index 2026.} Sydney: IEP.
\item
  International Federation of Robotics (2025). \emph{World Robotics 2025} (press release on 2024
  installations).
\item
  Karger, E., Rosenberg, J., Jacobs, Z., et al.~(2023). Forecasting existential risk: evidence from
  a long-run forecasting tournament. Forecasting Research Institute Working Paper 1.
\item
  Kulveit, J., Douglas, R., Ammann, N., Turan, D., Krueger, D., Duvenaud, D. (2025). Gradual
  disempowerment: systemic existential risks from incremental AI development. arXiv:2501.16946.
\item
  Kwa, T., West, B., et al.~(2025). Measuring AI ability to complete long tasks. METR.
  arXiv:2503.14499. With METR Time Horizon 1.1 data (May 2026).
\item
  Landay, K., Harms, P. D., Credé, M. (2019). Shall we serve the dark lords? A meta-analytic review
  of psychopathy and leadership. \emph{Journal of Applied Psychology} 104(1): 183-196.
  doi:10.1037/apl0000357
\item
  Milgram, S. (1974). \emph{Obedience to Authority.} Harper and Row. (Not re-checked.)
\item
  NOAA Climate Prediction Center (2026). ENSO diagnostic discussion, 10 September 2026.
\item
  Nord, M., et al.~(2026). \emph{Democracy Report 2026.} V-Dem Institute, University of Gothenburg.
\item
  O'Neill, B. C., et al.~(2017). The roads ahead: narratives for shared socioeconomic pathways
  describing world futures in the 21st century. \emph{Global Environmental Change} 42: 169-180. (Not
  re-checked.)
\item
  Ord, T. (2020). \emph{The Precipice: Existential Risk and the Future of Humanity.} Bloomsbury.
\item
  PRRI (2023). Findings from the 2023 American Values Survey. Washington DC: PRRI.
\item
  Ross, M. L. (2001). Does oil hinder democracy? \emph{World Politics} 53(3): 325-361. (Not
  re-checked.)
\item
  Rustad, S. A. (2025). \emph{Conflict Trends: A Global Overview, 1946-2024.} PRIO Paper. Oslo:
  PRIO.
\item
  Sanz-García, A., Gesteira, C., Sanz, J., García-Vera, M. P. (2021). Prevalence of psychopathy in
  the general adult population: a systematic review and meta-analysis. \emph{Frontiers in
  Psychology} 12: 661044. doi:10.3389/fpsyg.2021.661044
\item
  SIPRI (2025). Military Expenditure Database.
\item
  Stephan, M. J., Chenoweth, E. (2008). Why civil resistance works: the strategic logic of
  nonviolent conflict. \emph{International Security} 33(1): 7-44.
\item
  Straus, S. (2006). \emph{The Order of Genocide: Race, Power, and War in Rwanda.} Cornell
  University Press.
\item
  Sudduth, J. K. (2017). Strategic logic of elite purges in dictatorships. \emph{Comparative
  Political Studies} 50(13): 1768-1801.
\item
  Svolik, M. W. (2012). \emph{The Politics of Authoritarian Rule.} Cambridge University Press.
  doi:10.1017/CBO9781139176040
\item
  Tilly, C. (1990). \emph{Coercion, Capital, and European States, AD 990-1990.} Blackwell. (Not
  re-checked.)
\item
  Tsebelis, G. (2002). \emph{Veto Players: How Political Institutions Work.} Princeton University
  Press. (Not re-checked.)
\item
  UK Ministry of Defence (2025). Defence Intelligence update on Russian casualties, 14 October 2025.
\item
  UN Commission of Inquiry on Human Rights in the Democratic People's Republic of Korea (2014).
  Report A/HRC/25/63.
\item
  UN DESA (2024). \emph{World Population Prospects 2024.}
\item
  Valentino, B., Huth, P., Balch-Lindsay, D. (2004). ``Draining the sea'': mass killing and
  guerrilla warfare. \emph{International Organization} 58(2): 375-407. doi:10.1017/S0020818304582061
\item
  Weeks, J. L. P. (2012). Strongmen and straw men: authoritarian regimes and the initiation of
  international conflict. \emph{American Political Science Review} 106(2): 326-347.
  doi:10.1017/S0003055412000111
\item
  White House (2026). Presidential memorandum withdrawing the United States from 66 international
  organizations, 7 January 2026.
\end{itemize}

\end{document}
